\documentclass[11pt]{article}
\usepackage[T1]{fontenc}
\usepackage[utf8]{inputenc}
\usepackage{lmodern}
\usepackage[margin=1in]{geometry}
\usepackage{amsmath,amssymb,amsthm,mathtools}
\usepackage[shortlabels]{enumitem}
\usepackage{hyperref}
\newcommand{\OPHPaperReleaseID}{r1520}
\newcommand{\OPHPaperReleaseDate}{July 9, 2026}
\newcommand{\OPHPaperReleaseBanner}{%
\begin{center}
\small\textbf{Paper release:} \texttt{\OPHPaperReleaseID}\quad\textbf{Released:} \OPHPaperReleaseDate
\end{center}
}


\hypersetup{colorlinks=true,linkcolor=blue,citecolor=blue,urlcolor=blue}

\newtheorem{theorem}{Theorem}[section]
\newtheorem{corollary}[theorem]{Corollary}
\newtheorem{definition}[theorem]{Definition}
\newtheorem{assumption}[theorem]{Assumption}
\newtheorem{proposition}[theorem]{Proposition}
\newtheorem{claim}[theorem]{Claim}
\newtheorem{conjecture}[theorem]{Conjecture}
\theoremstyle{remark}
\newtheorem{remark}[theorem]{Remark}
\theoremstyle{plain}

\title{Reality as a Consensus Protocol:\\
The Fixed-Point Computation That Implements Physics}
\author{B.\ M\"uller \and Kai Xue \and Kale Arnav Anirudha}
\date{\OPHPaperReleaseDate}

\begin{document}
\maketitle
\OPHPaperReleaseBanner

\begin{abstract}
The consensus layer of OPH models a universe as a finite network of observer patches. Patches can carry mismatched overlap data. Repair moves reduce the mismatch until the public state settles. The computation is fixed-point selection of a public normal form, and cosmic history is the internal readout of that normal form.

The finite patch-net mechanics are the working layer used by the rest of the OPH stack. Accepted repair moves pass transactional validation and descend an exact measure until termination; atomic conflict-component commits and repair completeness give a schedule-independent normal form, and loop holonomy marks the exact places where local agreement fails globally. The same bookkeeping carries stable records with Born--L\"uders readout on the central event surface, refinement control, and a fair-block theorem for noisy consensus.
\end{abstract}

\section*{What This Paper Contributes}

The mathematics of repair is old in the right places. Constraint satisfaction, rewriting systems,
well-founded descent, local-diamond confluence, inverse limits, and recovery channels all enter this
paper with their standard meanings. OPH adds the physical reading: each variable is an observer
patch, each constraint is an overlap-visible record check, and each accepted rewrite is a local
repair move that must descend to the quotient seen by neighboring observers.

That turns consensus into the implementation layer of the theory. The paper proves when finite
patch repair terminates, when the terminal public state is schedule-independent, how boundary data
control uniqueness, and how cycle holonomy records the exact obstruction to global agreement. It
also explains why a bare overlap graph is only a finite constraint code. Stronger language such as
quantum error correction, min-cut distance, BFT liveness, or hardware speedup needs its own
certificate.

\section{Introduction}

This paper writes the OPH consensus picture in the simplest concrete form. A universe is represented as a finite graph of observer patches. Each patch carries local state data, neighboring patches compare those data on overlaps, and local repair moves try to reconcile any mismatch. The central mathematical question is whether this repair dynamics converges to one shared world and how the unavoidable obstructions are encoded when it does not.

The word ``dynamics'' is used here in the rewriting-system sense. The repair schedule selects
terminal normal forms and proves schedule independence under the stated hypotheses. Cosmic time is
the internal readout of the selected observer-facing structure. A conventional simulation computes
a surrogate history. OPH computes the fixed point whose internal record order is read as history.
The patch-net algorithm serves as a theorem device for fixed-point selection. The fundamental
description contains no global timeline on which spacetime contents are rendered.

A bare fixed point is not called an OPH simulation below.
Definition~\ref{def:oph-simulation} requires recovery-derived endogenous update,
nontrivial quotient-readable records, overlap repair with schedule-independent normal form,
elimination of a proper candidate basin in the selected boundary/sector fiber, and
implementation/clock closure. Theorem~\ref{thm:oph-simulation} proves those clauses on the
named finite branch and supplies a generic fixed-point and variational counterexample.

The resulting theorem package has two layers. The bare overlap network is only a finite
constraint code: its codewords are the globally overlap-consistent states. Quantum
error-correction, code-distance / min-cut resilience, exponential mixing, wall-clock BFT
liveness, and hardware speedup are separate certified branches with additional data. The first
core consensus layer concerns convergence: primitive recovery proposals are eligible only when
they satisfy the touched-overlap local-fit contract, but a proposal becomes physical only through
transactional acceptance with snapshot validation, protected-boundary preservation, and exact
well-founded descent. Connected conflict components commit atomically through canonical
union-collar payloads, so under repair completeness every initial quotient state has a unique
schedule-independent normal form. The stronger claim that all interiors with the same boundary
data settle to the same state requires preservation of that boundary data plus a unique consistent
extension in the corresponding fiber; the functional selected-fiber branch proves that condition
by constructing the single rooted extension and checking the remaining sector/holonomy
obstructions. The second concerns obstructions: pairwise overlap agreement does not ensure a
global solution, and the obstruction is holonomic. On the abelian branch it is the cycle sum of
edge data; on the genuinely noncentral branch it is a crossed-module \v{C}ech class.

Gauge symmetry enters as invariance under changes of hidden local representation that preserve
overlap data. When the repair step is read only on that overlap-invariant quotient, the
normal-form map descends to the gauge quotient, so physical uniqueness is a quotient statement.
On the quantum lift, the same quotient-local carrier determines a unique terminal state on every
declared physical observable algebra, even when microscopic representative lifts differ by gauge
or sector relabelings inside one quotient-local glued state.
The observation layer is carried by finite observer-accessible record algebras generated by
central or quantitatively stable approximately commuting projectors. These results form the
patch-net formulation used in the broader OPH literature, including the fixed-cutoff Bell/CHSH
package on the companion microphysics surface.

This paper proves a constraint-code firewall and nine core consensus results:

\begin{enumerate}
\item \textbf{Constraint-code firewall.} A finite overlap net defines a finite constraint code:
its codewords are exactly the globally consistent states, and \(C=\Phi^{-1}(0)\)
(Proposition~\ref{prop:constraint-code}). This is the default meaning of ``the overlap network is
a code.'' It is not, by itself, a quantum error-correcting code and does not imply a graph
min-cut formula for distance (Theorem~\ref{thm:no-free-mincut}).
\item \textbf{Asynchronous confluence.} For the declared accepted repair law, primitive recovery proposals pass through transactional snapshot/read/write validation, boundary/sector preservation, and exact descent. Atomic conflict-component commits give the local diamond on the physical quotient; under repair completeness, every fixed initial quotient state has a unique normal form, independent of update schedule (Theorem~\ref{thm:confluence}).
\item \textbf{Cycle obstruction.} For affine overlap constraints over an abelian group, global consistency holds if and only if the holonomy vanishes on every cycle (Theorem~\ref{thm:holonomy}). The parity triangle gives the minimal frustrated example, and Theorem~\ref{thm:higherdefects} extends the same logic to the crossed-module higher-gauge defect hierarchy used later in the framework.
\item \textbf{Gauge quotient, selected fibers, and observable-level confluence.} When local repair is induced on the overlap-invariant quotient, the normal-form map descends to that quotient, and the induced terminal state on every declared physical observable algebra is unique there even when microscopic representatives differ by gauge or sector relabelings inside one quotient-local glued state. If a boundary/sector map is preserved and each consistent boundary fiber has at most one quotient extension, then all initial states with that boundary value settle to the same quotient normal form; the functional selected-fiber branch proves a nontrivial case where multiple same-boundary interiors exist, inconsistent candidates are eliminated, and surviving candidates share one quotient normal form (Theorems~\ref{thm:gauge}, \ref{thm:boundary-conditioned-uniqueness}, \ref{thm:functional-selected-fiber}, Corollary~\ref{cor:branch-elimination}, and Theorem~\ref{thm:obsconfluence}).
\item \textbf{Refinement-limit consensus classes.} On a separated cofinal refinement system whose restriction maps commute with the finite-stage normal-form and holonomy maps, the quotient normal forms and holonomy obstructions assemble into unique inverse-limit classes with finite-stage visibility (Theorem~\ref{thm:refinement-consensus}).
\item \textbf{Coarse-graining compatibility.} On any refinement system whose coarse-graining maps shadow finite-stage normal forms and holonomy maps with declared errors, reconciling first and then coarse-graining gives the same macroscopic law data as coarse-graining first and then reconciling, up to those errors; in the exact natural case the error is zero (Theorem~\ref{thm:coarse-grain-consensus}).
\item \textbf{Record algebra and stability.} On the fixed-cutoff observer-accessible surface, central record projectors carry Born/L\"uders measurement directly, and approximate record projectors inherit explicit \((\varepsilon,\delta_{\mathrm{rec}})\) stability bounds on the same event surface (Theorem~\ref{thm:crdt}).
\item \textbf{OPH simulation firewall.} On a selected boundary/sector fiber,
``simulation'' requires recovery-derived endogenous update, nontrivial quotient-readable
records, strict overlap-repair descent with a schedule-independent normal form, collapse of a
proper candidate basin to one consistent quotient state, and implementation/clock closure.
Under the stated selected-fiber and record hypotheses, the finite OPH packet has this
certificate; a generic identity fixed point or constant variational functional does not
(Definition~\ref{def:oph-simulation}, Theorem~\ref{thm:oph-simulation}, and
Remark~\ref{rem:oph-simulation-falsifiers}).
\item \textbf{Distributed one-universe realization.} A worker implementation is a presentation
of one finite OPH universe only when the run starts from one global carrier and every distributed
event projects to a legal monolithic repair path, a physical stutter, or a certified rollback to
an earlier committed projection. Under those hypotheses, partition, worker count, schedule,
restart history, and repartition metadata do not change quotient observables or observer readouts
that factor through the monolithic normal form (Theorem~\ref{thm:distributed-one-universe} and
Corollary~\ref{cor:distributed-invariance}).
\item \textbf{Conditional noisy fair-block consensus.} If a noisy asynchronous implementation admits fair repair blocks, a uniform expected contraction toward the exact quotient normal-form set, and controlled within-block excursions, then all long runs stay in a controlled expected tube around that exact normal-form set. In singleton boundary/sector fibers, Lipschitz observer readouts are approximately schedule-independent, and a finite Markov-kernel certificate checks the block contraction on finite exported nets (Theorem~\ref{thm:noisy-fair-block-consensus}, Corollaries~\ref{cor:approx-schedule-independence}--\ref{cor:high-prob-noisy-tube}, and Proposition~\ref{prop:finite-fair-block-certificate}).
\end{enumerate}

The repair step itself is a concrete recovery move, not a free rewrite primitive. On the
fixed-cutoff collar branch, a local update is obtained from exact Markov splice or from a
declared Petz/Fawzi--Renner recovery channel and then read on overlap-invariant physical data.
The fixed-point theorem below isolates the separate branch conditions cleanly: the declared
repair law includes the touched-overlap local-fit contract for primitive proposals, while
transactional acceptance validates snapshots, protected data, unchanged registers, and exact
descent. The fixed-cutoff gluing package supplies the parenthesization-independent union-collar
payload used for atomic conflict-component commits. What sits above that step is repair
completeness and, on the Petz branch, the support/CPTP clause stated later in
Proposition~\ref{prop:petz-cptp}.
A nontrivial exported rooted-tree packet domain where these clauses are proved is recorded in
Definition~\ref{def:tree-packet-domain} and Theorem~\ref{thm:tree-packet-domain}.

We also define a fitness functional over a finite candidate space of reconciliation laws and prove
that replicator dynamics monotonically increases mean fitness (Theorem~\ref{thm:replicator}). This
gives a clean mathematical model for finite-candidate law selection, not a universality theorem or
a literal cosmological dynamics claim.

The results here are exact theorems about a computational model. The companion
OPH manuscript uses separate status categories for structural theorems, scaling
limits, quantitative particle outputs, and phenomenological
continuations~\cite{OPH}.


\section{Patch Nets, Overlaps, and Global Consistency}

\begin{definition}[Patch net]\label{def:patchnet}
Let $G=(V,E)$ be a finite connected graph. Each vertex $i\in V$ is an \textbf{observer patch} with finite local state space $S_i$. The global state space is
\[
\Sigma := \prod_{i\in V} S_i.
\]
For each edge $e=\{i,j\}\in E$, let $I_e$ be an interface alphabet and let
\[
\pi_{i,e}:S_i\to I_e,
\qquad
\pi_{j,e}:S_j\to I_e
\]
be the interface projection maps. A global state $s=(s_i)_{i\in V}\in\Sigma$ is \textbf{consistent on edge $e=\{i,j\}$} iff
\[
\pi_{i,e}(s_i)=\pi_{j,e}(s_j).
\]
The global consistency set is
\[
C:=\bigl\{s\in\Sigma:\forall\, e=\{i,j\}\in E,\ \pi_{i,e}(s_i)=\pi_{j,e}(s_j)\bigr\}.
\]
\end{definition}

For exposition we use a finite pairwise-overlap graph. The hypergraph version is straightforward: replace edges by hyperedges and pairwise equality by a common interface label on each hyperedge. Nothing in the proofs depends on the pairwise restriction.

The picture: each observer holds a local state, and neighboring observers share an interface through which they can compare notes. A universe-state is physically admissible exactly when all neighbors agree on their shared data. This is a constraint satisfaction problem (CSP), and the consistent states are the codewords.

\begin{definition}[Inconsistency potential]\label{def:potential}
For each edge $e$, choose a weight $w_e>0$ and a function $d_e:I_e\times I_e\to\mathbb{R}_{\ge 0}$ with $d_e(a,b)=0 \iff a=b$. On the declared fixed-cutoff branch, \(d_e\) is the overlap score used by the local acceptance contract on that interface. Define
\[
\Phi(s)
:=
\sum_{e=\{i,j\}\in E}
w_e\, d_e\!\bigl(\pi_{i,e}(s_i),\pi_{j,e}(s_j)\bigr).
\]
Then $s\in C \iff \Phi(s)=0$.
\end{definition}

\begin{proposition}[Bare overlap nets are finite constraint codes]\label{prop:constraint-code}
For every finite patch net of Definition~\ref{def:patchnet}, the consistency set \(C\subseteq
\Sigma\) is a finite constraint code whose codewords are exactly the globally overlap-consistent
states. With the mismatch potential of Definition~\ref{def:potential},
\[
C=\Phi^{-1}(0).
\]
This proposition is the theorem-grade content of the unqualified phrase ``the overlap network is
a code.'' It supplies a finite constraint code, not a quantum error-correcting code, not a
topological code, and not a graph-theoretic formula for code distance.
\end{proposition}

\begin{proof}
Finiteness follows from \(\Sigma=\prod_i S_i\) with finite \(S_i\). The definition of \(C\) is a
finite family of interface-equality constraints, so \(C\) is the set of satisfying assignments.
Since each \(w_e>0\) and \(d_e(a,b)=0\) exactly when \(a=b\), every term in \(\Phi\) is
nonnegative and vanishes exactly on a satisfied edge constraint. Therefore all edge constraints
hold if and only if \(\Phi(s)=0\).
\end{proof}

So $\Phi$ is the total disagreement energy of the universe. Consistent states have zero energy. Everything else is frustrated.


\section{Asynchronous Reconciliation and the Main Fixed-Point Theorem}

\begin{definition}[Recovery-derived local repair law]\label{def:repair}
Fix for each patch \(i\) a finite collar chart \(A_i\!-\!B_i\!-\!D_i\) around the overlaps touched
by \(i\), together with a fixed local decoder from repaired collar data back to the finite patch
label at \(i\). A \textbf{law} \(\lambda\) is a family of local repair maps
\[
T_i^\lambda:\Sigma\to\Sigma
\qquad (i\in V)
\]
such that \(T_i^\lambda\) changes only the state of patch \(i\) (or, more generally, only a
bounded neighborhood of \(i\)), and the local update is induced by one of the declared OPH
recovery moves on that collar, where \(\omega_{A_iB_i}(s)\) denotes the \(A_i\cup B_i\) marginal
of the input collar state encoded by \(s\):
\begin{enumerate}[label=(\roman*),leftmargin=2em]
\item exact Markov splice on the collar, using Theorem~\ref{thm:splice} when
      \(I(A_i:D_i\mid B_i)=0\); or
\item a declared recoverability channel
\[
(\mathrm{id}_{A_i}\otimes \mathcal R_i)(\omega_{A_iB_i}(s)),
\qquad
\mathcal R_i=\mathcal R_{\sigma_i,\mathcal N_i},
\]
      with \(\mathcal R_i\) in the Petz/Fawzi--Renner class of
      Definition~\ref{def:oph-repair-map} and Theorem~\ref{thm:splice}.
\end{enumerate}
The decoder back to \(S_i\) is bookkeeping for the finite patch presentation; the physical content
is the repaired collar state on the declared fixed-cutoff branch. Write
\(s\rightsquigarrow_i t\) iff \(t=T_i^\lambda(s)\neq s\). These are primitive proposals, not yet
accepted physical rewrite steps. A proposal is committed only through the transactional acceptance
layer below.
\end{definition}

\begin{definition}[Touched-overlap potential and accepted local-fit contract]\label{def:localfit}
For each repair site \(i\), let
\[
E_i^{\mathrm{touch}}
:=
\bigl\{
e\in E:\text{the interface data on }e\text{ may change under }T_i^\lambda
\bigr\}.
\]
Define the touched-overlap potential
\[
\Phi_i(s)
:=
\sum_{e=\{u,v\}\in E_i^{\mathrm{touch}}}
w_e\, d_e\!\bigl(\pi_{u,e}(s_u),\pi_{v,e}(s_v)\bigr).
\]
On the declared fixed-cutoff branch, a recovery-derived primitive candidate is eligible for
transactional commitment only if it strictly lowers this touched-overlap score:
\[
s\rightsquigarrow_i t \implies \Phi_i(t)<\Phi_i(s).
\]
This is the patch-net form of the regulator-side monotone local-fit contract carried by the
declared repair package.
\end{definition}

\begin{definition}[Overlap-associative union-collar gluing]\label{def:assoc}
Fix \(s\in\Sigma\) and two enabled primitive proposals \(s\rightsquigarrow_i t\),
\(s\rightsquigarrow_j u\). Write
\[
E_{ij}^{\mathrm{touch}}
:=
E_i^{\mathrm{touch}}\cup E_j^{\mathrm{touch}}.
\]
The declared repair branch is \textbf{overlap-associative} if the following hold.
\begin{enumerate}[label=(\roman*),leftmargin=2em]
\item If \(E_i^{\mathrm{touch}}\cap E_j^{\mathrm{touch}}=\varnothing\), the two local proposals
have disjoint support on the declared branch and therefore commute if they are committed as
separate transactions.
\item If \(E_i^{\mathrm{touch}}\cap E_j^{\mathrm{touch}}\neq\varnothing\), there is a finite union
collar \(U_{ij}\) covering the interfaces in \(E_{ij}^{\mathrm{touch}}\) such that the physical
glued state on \(U_{ij}\) is parenthesization-independent on the quotient, in the sense of
Proposition~\ref{prop:assocglue}, and the local decoders/lifts of
Definition~\ref{def:repair} are restriction-compatible on nested collars.
\end{enumerate}
This is the concrete compatibility package used below to build canonical aggregate payloads for
conflicting repair components.
\end{definition}

\begin{definition}[Transactional acceptance layer]\label{def:transactional-repair}
Let \(X\) denote the finite physical presentation on which a repair step is being checked: before
quotienting one may take \(X=\Sigma\), and on the physical branch \(X\) is the quotient by hidden
representatives. Registers are the finite patch, interface, sector, and record coordinates. Fix a
boundary/sector/holonomy record map
\[
B:X\to\mathcal B
\]
and a well-founded exact measure
\[
\mu:X\to(W,\prec),
\]
for example a lexicographic integer vector
\((N_{\mathrm{hard}},\Phi,N_{\mathrm{unresolved}})\).

A prepared transaction is a tuple
\[
\tau=(R_\tau,W_\tau,\sigma_\tau,p_\tau)
\]
with read set, write set, read snapshot, and payload. It commits at \(x\in X\) only if all of the
following hold:
\begin{enumerate}[label=(\roman*),leftmargin=2em]
\item the snapshot is current: \(x|_{R_\tau}=\sigma_\tau\);
\item the payload changes no register outside \(W_\tau\);
\item boundary, sector, holonomy, and protected record data are preserved:
      \(B(\operatorname{Apply}_\tau(x))=B(x)\);
\item exact descent holds:
      \[
      \mu(\operatorname{Apply}_\tau(x))\prec \mu(x).
      \]
\end{enumerate}
A stale, aborted, ambiguous, or obstructed transaction is not a rewrite step. The read set is
\textbf{validation-complete} for \(\mu\): if a term of \(\mu\) can change when a register in
\(W_\tau\) changes, the whole support of that term lies in \(R_\tau\). For seam potentials this
means reading both endpoints of every seam whose score can change.

At a state \(x\), form the conflict graph of enabled primitive repair proposals, with
\[
\tau\#\sigma
\quad\Longleftrightarrow\quad
W_\tau\cap(R_\sigma\cup W_\sigma)\ne\varnothing
\ \text{or}\
W_\sigma\cap(R_\tau\cup W_\tau)\ne\varnothing .
\]
Each connected conflict component \(K\) is replaced by exactly one canonical aggregate transaction
\(\tau_K\), computed on the union-collar or conflict-component support. Primitive members of
\(K\) never commit separately. Write \(x\to y\) for a successful aggregate transaction commit, and
let \(\to^*\) be its reflexive-transitive closure. A state is a \textbf{normal form} when no
aggregate transaction is enabled.
\end{definition}

\begin{proposition}[Validation support for local mismatch measures]\label{prop:validation-support}
Suppose the exact repair measure has finite local supports,
\[
\mu(x)=\sum_{a\in A}\mu_a(x|_{S_a}).
\]
If a transaction writes \(W\), then only terms with \(S_a\cap W\ne\varnothing\) can change.
Consequently, the descent validation is snapshot-local once the read set contains
\[
R^\mu(W):=\bigcup_{a:S_a\cap W\ne\varnothing}S_a .
\]
For the OPH overlap potential
\[
\Phi(x)=\sum_{e=\{i,j\}}w_e\,d_e\!\bigl(\pi_{i,e}(x_i),\pi_{j,e}(x_j)\bigr),
\]
this means reading both endpoints of every overlap whose score may change when a written register
changes.
\end{proposition}

\begin{proof}
If \(S_a\cap W=\varnothing\), the transaction leaves every register read by \(\mu_a\) unchanged,
so that term cancels between the pre- and post-state. Every remaining term is determined by the
payload together with the restriction to \(R^\mu(W)\). The displayed edge-potential formula has
one two-endpoint support for each overlap term, giving the stated seam rule.
\end{proof}

\begin{remark}[Inputs and branch conditions]\label{rem:repair-burdens}
The repair step is therefore not an abstract rewrite primitive. Its declared inputs are the
fixed-cutoff collar chart, either exact Markov splice or a chosen Petz/Fawzi--Renner recovery
channel, a local decoder/lift back to the finite patch presentation, and the touched-overlap
local-fit contract of Definition~\ref{def:localfit}, together with the support-local
disjoint-commutation clause and the restriction-compatible union-collar package of
Definition~\ref{def:assoc}. The
parenthesization-independent quotient-local glue used there is supplied by
Proposition~\ref{prop:assocglue} from the fixed-cutoff center-sector / higher-gauge gluing
package. The actual accepted step is the validated aggregate transaction of
Definition~\ref{def:transactional-repair}. The theorem package takes repair completeness as an
explicit branch condition. On the Petz branch, full CPTP action on all inputs also requires the support clause
recorded in Proposition~\ref{prop:petz-cptp}. On broader branches one must also prove that
the declared union-collar compatibility is preserved under refinement or branch change.
\end{remark}

The theorem package separates the imported repair-law data from the theorem-local inputs cleanly:

\begin{assumption}[Repair completeness]\label{ass:complete}
For the accepted transactional relation \(\to\), \(s\in C\) if and only if no aggregate repair
transaction is enabled at \(s\).
\end{assumption}

Normal forms are exactly the globally consistent states. The dynamics is neither too weak (missing some inconsistencies) nor too strong (repairing things that were fine).

\begin{definition}[Verified rooted-tree packet-net domain]\label{def:tree-packet-domain}
Fix a finite rooted tree \(T=(V,E,r)\). For each non-root vertex \(i\), write \(p(i)\) for its
parent and write \(w_i>0\) for the weight of the edge \(\{p(i),i\}\). Let \(A\) be a finite
packet alphabet with \(|A|\ge2\), and let \(K_i\) be a finite hidden-label set. The patch state
space is
\[
S_i=A\times K_i,
\qquad
s_i=(x_i,k_i).
\]
For every edge \(e=\{i,j\}\), the interface alphabet is \(I_e=A\), and both endpoint projections
read the packet component:
\[
\pi_{i,e}(x_i,k_i)=x_i,
\qquad
\pi_{j,e}(x_j,k_j)=x_j.
\]
Thus \(e\) is consistent exactly when \(x_i=x_j\). Choose the weights so that
\[
w_i>\sum_{j:p(j)=i}w_j
\qquad
\text{for every non-root } i,
\]
with an empty sum equal to \(0\). Define
\[
\Phi(s)=\sum_{\{p(i),i\}\in E}w_i\,\mathbf 1[x_i\ne x_{p(i)}].
\]
The repair map at a non-root vertex \(i\) is
\[
T_i(s)_i=(x_{p(i)},k_i),
\]
with all other vertices unchanged; if \(x_i=x_{p(i)}\), the map is a no-op. The root repair map is
the identity. The hidden labels \(k_i\) are acted on by arbitrary finite gauge relabelings and are
not read by any interface projection.
\end{definition}

\begin{theorem}[Rooted-tree packet repair completeness and quotient closure]\label{thm:tree-packet-domain}
On the domain of Definition~\ref{def:tree-packet-domain}, Assumption~\ref{ass:complete} is a
theorem. Every enabled repair strictly decreases \(\Phi\). Every maximal asynchronous repair run
terminates at the unique state
\[
x_i=x_r\quad\text{for all }i\in V,
\]
with all hidden labels \(k_i\) unchanged. The normal-form map descends to the quotient by hidden
gauge relabeling. If several tree repairs lie in one conflict component, the canonical aggregate
transaction applies them in increasing tree depth and is the same as the corresponding serial
tree repair path. The four-vertex instance with edges \(r\!-\!a\), \(a\!-\!b\), \(a\!-\!c\),
alphabet \(A=\mathbb Z_3\), hidden labels \(K_i=\mathbb Z_2\), and weights \(5,1,1\) is exported
as \path{code/consensus/runs/verified_tree_packet_net_domain.json}.
\end{theorem}

\begin{proof}
First, repair completeness is immediate from the rooted tree. If \(s\in C\), every edge has
\(x_i=x_{p(i)}\), so every non-root repair is a no-op. Conversely, if every repair is a no-op,
then \(x_i=x_{p(i)}\) for every non-root vertex, hence every edge is consistent and \(s\in C\).

Let \(i\ne r\) be enabled. Only the parent edge \(\{p(i),i\}\) and child edges
\(\{i,j\}\) with \(p(j)=i\) can change their contribution to \(\Phi\). The parent edge changes
from inconsistent to consistent, contributing \(-w_i\). Each child edge can increase by at most
its weight. Therefore
\[
\Phi(T_i(s))-\Phi(s)
\le
-w_i+\sum_{j:p(j)=i}w_j
<0.
\]
So every enabled repair is accepted by the Lyapunov contract.

Since \(\Phi\) takes finitely many values, every maximal repair run terminates. At a terminal
state no non-root vertex differs from its parent, so the terminal packet label is \(x_r\) on every
vertex. The root label and every hidden label are invariant under all repairs, so the terminal
state is unique and independent of update order. Because interface projections, enabledness,
\(\Phi\), and the repair maps depend only on \(x_i\), arbitrary relabelings of the hidden
\(K_i\) commute with quotienting:
\[
q\circ T_i=\overline T_i\circ q.
\]
Thus the normal-form map descends to the hidden-label quotient.
\end{proof}

\begin{corollary}[Physical-law map on the verified packet domain]\label{cor:tree-packet-law}
On the rooted-tree packet domain, every gauge-invariant observable
\[
M:\prod_i(A\times K_i)\to Y
\]
has a schedule-independent repaired value
\[
M(\operatorname{nf}(s)),
\]
and this value depends only on the quotient class of \(s\). Hence the promoted normal-form map is
usable as physical law on this verified packet branch without adding a representative-level
gauge-covariance assumption.
\end{corollary}

\begin{proof}
Theorem~\ref{thm:tree-packet-domain} gives a unique terminal state for every asynchronous repair
schedule and shows that the normal-form map descends to the quotient by hidden-label relabeling.
A gauge-invariant observable factors through that quotient, so its value on the terminal state is
independent of both the repair schedule and the hidden representative.
\end{proof}

\begin{proposition}[Classical full-support Petz packet domain]\label{prop:classical-petz-packet}
Let \(B\) and \(D\) be finite packet alphabets, let the collar algebras be diagonal, and let
\(\mathcal N:\mathbb C^{B\times D}\to\mathbb C^B\) be the marginal channel
\((\mathcal N p)(b)=\sum_d p(b,d)\). Fix a reference state \(\sigma_{BD}\) with
\(\sigma_B(b)>0\) for every \(b\). Let
\[
\gamma_\sigma:=\min_{b\in B}\sigma_B(b)>0.
\]
Define the Petz recovery channel
\[
\mathcal R_{\sigma,\mathcal N}(\mu)(b,d)
=
\mu(b)\,\sigma(d\mid b),
\qquad
\sigma(d\mid b):=\frac{\sigma_{BD}(b,d)}{\sigma_B(b)}.
\]
Then \(\mathcal R_{\sigma,\mathcal N}\) is stochastic, completely positive and trace preserving
on the diagonal algebra, and \(\ell^1\)-contractive:
\[
\|\mathcal R_{\sigma,\mathcal N}(\mu)-\mathcal R_{\sigma,\mathcal N}(\nu)\|_1
\le
\|\mu-\nu\|_1.
\]
The support inverse is uniformly bounded on this domain by
\[
\|\sigma_B^{-1/2}\|\le \gamma_\sigma^{-1/2}.
\]
If a collar state has the exact classical Markov form
\[
\omega_{ABD}(a,b,d)=\omega_{AB}(a,b)\sigma(d\mid b),
\]
then \((\mathrm{id}_A\otimes\mathcal R_{\sigma,\mathcal N})(\omega_{AB})=\omega_{ABD}\).
The support obstruction is exact: if \(\sigma_B(b)=0\) and an input assigns mass to \(b\), the
Petz inverse on that sector is undefined unless the channel domain is restricted or a separate
trace-preserving completion is declared.
\end{proposition}

\begin{proof}
The displayed formula is the finite diagonal specialization of the Petz map
\[
\sigma_{BD}^{1/2}
\left(\sigma_B^{-1/2}\mu\,\sigma_B^{-1/2}\otimes \mathbf 1_D\right)
\sigma_{BD}^{1/2}.
\]
Full support of \(\sigma_B\) makes the inverse well defined, with support gap
\(\gamma_\sigma>0\), hence \(\|\sigma_B^{-1/2}\|\le\gamma_\sigma^{-1/2}\). Nonnegativity is immediate, and
\[
\sum_{b,d}\mathcal R_{\sigma,\mathcal N}(\mu)(b,d)
=
\sum_b \mu(b)\sum_d\sigma(d\mid b)
=
\sum_b\mu(b),
\]
so the map is trace preserving. Diagonal positive trace-preserving maps are completely positive on
the diagonal algebra. For signed diagonal inputs,
\[
\|\mathcal R_{\sigma,\mathcal N}(\mu)-\mathcal R_{\sigma,\mathcal N}(\nu)\|_1
=
\sum_{b,d}|\mu(b)-\nu(b)|\,\sigma(d\mid b)
=
\sum_b|\mu(b)-\nu(b)|,
\]
which gives the stated contraction. The exact Markov recovery identity follows by substitution.
If \(\sigma_B(b)=0\), the factor \(\sigma_B^{-1/2}\) has no value on that sector; mass placed
there by an input is outside the Petz support domain. That is the claimed obstruction.
\end{proof}

\begin{proposition}[Accepted repair moves are decreasing]\label{prop:lyapunov-descent}
For the accepted transactional repair law of
Definitions~\ref{def:repair}--\ref{def:transactional-repair}, every enabled repair strictly
decreases the declared exact measure:
\[
s\to t \implies \mu(t)\prec\mu(s).
\]
On the scalar \(\Phi\)-branch, this specializes to \(\Phi(t)<\Phi(s)\).
\end{proposition}

\begin{proof}
This is part of the commit validation in Definition~\ref{def:transactional-repair}. For a
single-site scalar-\(\Phi\) transaction it reduces to the touched-overlap contract of
Definition~\ref{def:localfit}; for an aggregate conflict component it is checked on the aggregate
payload before the component can commit.
\end{proof}

\begin{proposition}[Termination from the OPH Lyapunov functional]\label{prop:lyapunov-termination}
Under Proposition~\ref{prop:lyapunov-descent}, every repair sequence is finite; equivalently, the
repair relation \(\to\) is terminating.
\end{proposition}

\begin{proof}
Along any nontrivial repair step \(s\to t\), Proposition~\ref{prop:lyapunov-descent} gives
\(\mu(t)\prec\mu(s)\). The measure order is well founded, so an infinite strictly descending chain
is impossible. On the scalar-\(\Phi\) branch this is the older finite-value argument on
\(\Phi(\Sigma)\).
\end{proof}

\begin{proposition}[Finite repair step bound]\label{prop:finite-step-bound}
Let \(s_0\in\Sigma\). Under Proposition~\ref{prop:lyapunov-descent}, every accepted repair run
starting at \(s_0\) has length at most the number of distinct \(\mu\)-values reachable below
\(\mu(s_0)\) minus one. On the scalar-\(\Phi\) branch this gives
\[
\bigl|\{\,\Phi(s):s\in\Sigma,\ \Phi(s)\le \Phi(s_0)\,\}\bigr|-1
\le |\Phi(\Sigma)|-1.
\]
If, additionally, \(\Phi\) is integer-valued, or more generally every accepted move lowers
\(\Phi\) by at least a fixed \(\eta>0\), then the run length satisfies
\[
T(s_0)\le \left\lceil\frac{\Phi(s_0)}{\eta}\right\rceil .
\]
This is a finite descent bound in repair steps only. It is not a wall-clock bound, not a
probability-one scheduling statement, and not spectral or exponential convergence.
\end{proposition}

\begin{proof}
The values of \(\mu\) strictly decrease along the run, so no \(\mu\)-value can occur twice. This
gives the finite reachable-value bound. On the scalar-\(\Phi\) branch, the values of \(\Phi\)
strictly decrease along the run, so no value in
\(\{\,\Phi(s):s\in\Sigma,\ \Phi(s)\le \Phi(s_0)\,\}\) can occur twice. This gives the finite
value-set bound. If each move lowers \(\Phi\) by at least \(\eta\), after \(T\) moves the value has
dropped by at least \(T\eta\). Since \(\Phi\ge0\), \(T\eta\le\Phi(s_0)\), giving the displayed
ceiling bound.
\end{proof}

\begin{remark}[Termination is not confluence]\label{rem:termination-not-confluence}
Proposition~\ref{prop:lyapunov-termination} proves only that accepted repair runs stop. It does not
prove that two update orders stop at the same physical state. Schedule-independence enters only after
the local-diamond property of Proposition~\ref{prop:local-confluence} is combined with termination
via Newman's lemma and with repair completeness in Assumption~\ref{ass:complete}. If two accepted
repair schedules from the same initial state reach different observer-facing quotient normal forms,
with no declared holonomy or higher-gauge obstruction and no mere gauge-representative difference,
then the proposed repair law is not OPH-admissible as a consensus mechanism.
\end{remark}

\begin{proposition}[Local confluence from atomic conflict-component commits]\label{prop:local-confluence}
For the accepted repair law of
Definitions~\ref{def:repair}--\ref{def:transactional-repair}, the transactional repair relation
is locally confluent.
\end{proposition}

\begin{proof}
Take a one-step peak \(t\leftarrow s\to u\), and let the two committed aggregate transactions
come from enabled conflict components \(K\) and \(L\) at \(s\). If \(K=L\), the canonical
component rule in Definition~\ref{def:transactional-repair} gives the same aggregate payload, so
\(t=u\).

Assume \(K\ne L\). Then no write set in one component intersects any read-or-write set in the
other. Therefore each transaction's snapshot is still current after the other commits, and the
unchanged-outside-write test is unaffected. Validation-complete read sets ensure that every
measure term whose value can change under one component was read by that component, so the
descent check made at \(s\) remains the same after the other disjoint component is committed.
Boundary and sector preservation are likewise checked on disjoint protected supports or on
unchanged data. Thus each transaction can be committed after the other, and the two applications
write disjoint registers to the same values. The two length-two paths therefore meet at an
identical state \(v\). This is the local diamond.
\end{proof}

\begin{theorem}[Asynchronous confluence / fixed-point law]\label{thm:confluence}
For the accepted repair law of Definitions~\ref{def:repair}--\ref{def:transactional-repair},
under Assumption~\ref{ass:complete}, each fixed initial state $s\in\Sigma$ has a unique normal form
\[
\operatorname{nf}_\lambda(s)\in C,
\]
and every maximal asynchronous repair execution from $s$ terminates at that same state. The terminal state is independent of update order because descent, local confluence on the physical quotient, and repair completeness all hold; termination alone is not enough.
\end{theorem}

\begin{proof}
By Proposition~\ref{prop:lyapunov-termination}, the repair relation \(\to\) is terminating.
By Proposition~\ref{prop:local-confluence}, it is locally confluent. Newman's
lemma~\cite{newman42} therefore makes it confluent. A terminating confluent repair relation has a
unique normal form reachable from each fixed initial state. By Assumption~\ref{ass:complete}, the
normal forms are exactly \(C\). Every maximal execution terminates and reaches
\(\operatorname{nf}_\lambda(s)\), and that terminal state is independent of update order.
\end{proof}

\begin{corollary}[Objective law is schedule-independent]\label{cor:objective}
Let $M:\Sigma\to Y$ be any observable. Under the hypotheses of Theorem~\ref{thm:confluence}, $M(\operatorname{nf}_\lambda(s))$ is independent of the asynchronous update schedule. If physical law is identified with the map $s\mapsto M(\operatorname{nf}_\lambda(s))$, then physical law is objective.
\end{corollary}

\begin{proof}
All schedules from the same initial $s$ terminate at $\operatorname{nf}_\lambda(s)$, so all yield the same $M$-value.
\end{proof}

Here objectivity is identified with schedule-independent convergence of the repair
dynamics; Theorem~\ref{thm:obsconfluence} sharpens this to schedule-independent convergence of the
physical observable algebra even when microscopic representatives are not unique.

\begin{remark}[Quantifier on normal-form uniqueness]\label{rem:normal-form-quantifier}
Theorem~\ref{thm:confluence} proves uniqueness from a fixed initial state, and
Theorem~\ref{thm:gauge} below turns that into uniqueness from a fixed physical quotient state. It
does not say that all possible initial data settle to one universal state. Different boundary
conditions, conserved charges, root packets, holonomy sectors, or external records can
legitimately determine different normal forms.
\end{remark}


\section{Why Local Agreement Is Not Enough: Cycle Holonomy and Frustration}

Pairwise neighbor agreement does not by itself imply global consistency.

\begin{theorem}[Cycle-obstruction / holonomy criterion]\label{thm:holonomy}
Let $A$ be an abelian group, and let $G=(V,E)$ be a connected graph with an arbitrary orientation on each edge. For each oriented edge $e:u\to v$, assign a label $b_e\in A$. Consider the affine consistency equations
\[
x_v - x_u = b_e
\qquad
\text{for every oriented edge } e:u\to v,
\]
where $x_v\in A$ are unknown patch labels. A global solution $x:V\to A$ exists if and only if for every cycle $C\subseteq G$,
\[
\sum_{e\in C}\varepsilon_C(e)\, b_e = 0,
\]
where $\varepsilon_C(e)=+1$ if the cycle traverses $e$ in the chosen orientation and $-1$ otherwise.
\end{theorem}

\begin{proof}
\textbf{Necessity.} Suppose $x$ is a solution. Summing the edge equations around any cycle $C$,
\[
\sum_{e:u\to v\in C}\varepsilon_C(e)\,(x_v-x_u)
=
\sum_{e\in C}\varepsilon_C(e)\,b_e.
\]
The left side telescopes to $0$ because every vertex appears once with $+$ sign and once with $-$ sign.

\textbf{Sufficiency.} Fix a root $r\in V$. For any vertex $v$, choose a path $P_{r\to v}$ and define
\[
x_v := \sum_{e\in P_{r\to v}} \varepsilon_{P_{r\to v}}(e)\, b_e,
\qquad x_r:=0.
\]
If $P$ and $P'$ are two paths from $r$ to $v$, traversing $P$ followed by the reverse of $P'$ yields a cycle. By the vanishing-holonomy assumption, the total signed sum is zero, so $x_v$ is well-defined. For any edge $e:u\to v$, extending a path to $u$ by that edge gives $x_v = x_u + b_e$.
\end{proof}

\begin{corollary}[Parity triangle: pairwise consistency is not enough]\label{cor:triangle}
Take $A=\mathbb{Z}_2$ on the triangle $A$-$B$-$C$-$A$ with edge labels $b_{AB}=0$, $b_{BC}=0$, $b_{CA}=1$. Each individual edge equation is satisfiable. But the global system is not: the cycle sum is $0\oplus 0\oplus 1 = 1\neq 0$.
\end{corollary}

This example shows that pairwise consistency does not imply a global solution. The obstruction is
carried by the cycle.

\begin{corollary}[Stable defects as frustrated holonomy]\label{cor:defects}
Define the defect energy
\[
\Phi_b(x) := \sum_{e:u\to v\in E} w_e\,\mathbf{1}\!\left[x_v - x_u \neq b_e\right].
\]
If the cycle-holonomy condition fails, then $\min_{x:V\to A}\Phi_b(x)>0$. Every minimizer contains irreducible residual inconsistency.
\end{corollary}

\begin{proof}
If $\min_x\Phi_b(x)=0$, some assignment satisfies all edge equations, contradicting Theorem~\ref{thm:holonomy}.
\end{proof}

Residual inconsistencies of this type cannot be removed by local repair moves. In the OPH
interpretation they are stable topological defects of the reconciliation dynamics.

\begin{theorem}[Higher-gauge defect hierarchy]\label{thm:higherdefects}
Let a finite overlap nerve carry crossed-module defect data \((g_{ij},h_{ijk})\) for a compact crossed module
\[
H\xrightarrow{\partial} G.
\]
Under local rechartings by
\[
C^1(N,H)\rtimes C^0(N,G),
\]
the nonabelian \v{C}ech class
\[
q=[(g,h)]\in \check H^2(N,H\to G)
\]
is invariant. Strict global reconciliation exists if and only if \(q=0\), and nonzero \(q\) labels stable fixed-cutoff higher-gauge defects.
\end{theorem}

\begin{proof}
The allowed rechartings are exactly the crossed-module coboundaries. They preserve the class and
strictify the weak gluing data precisely when the class vanishes. A nonzero class therefore gives
a topologically protected residual obstruction in the full crossed-module defect hierarchy,
including the abelian truncation as its first layer.
\end{proof}


\section{Gauge Symmetry as Implementation Hiding}

\begin{definition}[Gauge action]\label{def:gauge}
Let $\Gamma = \prod_{i\in V}\Gamma_i$ act on $\Sigma$ componentwise. The action is a \textbf{gauge action} if for every $e=\{i,j\}\in E$,
\[
\pi_{i,e}(\gamma_i\cdot x)=\pi_{i,e}(x)
\qquad
\forall\, x\in S_i,\ \forall\, \gamma_i\in \Gamma_i.
\]
Gauge changes alter hidden local representations but do not alter overlap data.
\end{definition}

Gauge transformations change hidden local representations while leaving overlap data fixed. Write
\[
q:\Sigma\to \Sigma/\Gamma,
\qquad
q(s)=[s],
\]
for the gauge-orbit map. A \textbf{physical repair law} is the family of quotient-local maps
\[
\overline T_i^\lambda:\Sigma/\Gamma\to \Sigma/\Gamma
\]
induced by the recovery-derived collar updates of Definition~\ref{def:repair}. A
\textbf{representative repair family} is any choice of lifts
\[
T_i^\lambda:\Sigma\to\Sigma
\]
such that
\[
q\circ T_i^\lambda
=
\overline T_i^\lambda\circ q.
\]
This is the finite patch-net form of saying that the repair step is defined first on
overlap-invariant physical data and only then lifted to hidden representatives. The rooted-tree
packet domain of Theorem~\ref{thm:tree-packet-domain} proves repair completeness and quotient
descent on a nontrivial exported packet net, while Proposition~\ref{prop:classical-petz-packet}
proves the classical full-support Petz clause used by that domain. A broader fixed-cutoff
branch requires repair completeness for the chosen exported packet net and the
support/CPTP clause on the Petz branch where that channel is used. The touched-overlap local-fit
contract gives scalar \(\Phi\)-descent for primitive proposals on the corresponding branch, while
the transaction layer supplies exact accepted-step descent. Proposition~\ref{prop:assocglue}
together with Definition~\ref{def:assoc} supplies the quotient-local compatibility package used
to build canonical aggregate payloads; the local diamond then comes from atomic conflict-component
commits. Stability of that package under refinement or branch change is a separate question.
The point here is also that no extra gauge-covariance axiom is needed once repair is formulated
on the quotient.

\begin{theorem}[Gauge quotient theorem]\label{thm:gauge}
Under the gauge action of Definition~\ref{def:gauge}, any representative lift of a physical repair
law as just defined, and the hypotheses of Theorem~\ref{thm:confluence},
\[
q\bigl(\operatorname{nf}_\lambda(\gamma\cdot s)\bigr)
=
q\bigl(\operatorname{nf}_\lambda(s)\bigr)
\qquad
\forall\, \gamma\in\Gamma,\ \forall\, s\in\Sigma.
\]
Hence the normal-form map descends to the quotient:
\[
\overline{\operatorname{nf}}_\lambda:\Sigma/\Gamma \to \Sigma/\Gamma,
\qquad
[s]\mapsto [\operatorname{nf}_\lambda(s)].
\]
\end{theorem}

\begin{proof}
Let \(s\to t\) be an accepted aggregate transaction, and let \(\overline\tau\) be its quotient
payload. If \(s'=\gamma\cdot s\), the representative-lift hypothesis gives an accepted lift
\(s'\to t'\) of the same quotient transaction, and
\[
q(t')
=
\overline\tau\bigl(q(s')\bigr)
=
\overline\tau\bigl(q(s)\bigr)
=
q(t).
\]
Thus gauge-equivalent inputs induce the same repaired orbit, and by induction the orbit reached
after any repair sequence depends only on the initial orbit. Every maximal repair sequence from
$s$ ends at $\operatorname{nf}_\lambda(s)$ by Theorem~\ref{thm:confluence}, so the terminal orbit
$q(\operatorname{nf}_\lambda(s))$ depends only on $q(s)=[s]$. This makes
\[
[s]\longmapsto [\operatorname{nf}_\lambda(s)]
\]
well-defined on $\Sigma/\Gamma$.
\end{proof}

\begin{theorem}[Boundary-conditioned quotient uniqueness]\label{thm:boundary-conditioned-uniqueness}
Let
\[
B:\Sigma/\Gamma\to\mathcal B
\]
record fixed external boundary data, conserved charge, root packet, holonomy sector, or task
input. Assume accepted quotient repairs preserve \(B\):
\[
[s]\to [t]\implies B([s])=B([t]).
\]
For \(b\in\mathcal B\), define the consistent quotient fiber
\[
C_b:=\{\,x\in q(C):B(x)=b\,\}.
\]
If every \(C_b\) has at most one element, then all initial states with the same boundary value
settle to the same observer-facing quotient normal form:
\[
B([s])=B([s'])\implies
\overline{\operatorname{nf}}_\lambda([s])
=
\overline{\operatorname{nf}}_\lambda([s']).
\]
\end{theorem}

\begin{proof}
Theorem~\ref{thm:gauge} gives unique quotient normal forms
\(\overline{\operatorname{nf}}_\lambda([s])\) and
\(\overline{\operatorname{nf}}_\lambda([s'])\) in \(q(C)\). Boundary preservation along accepted
repair sequences gives
\[
B(\overline{\operatorname{nf}}_\lambda([s]))=B([s]),
\qquad
B(\overline{\operatorname{nf}}_\lambda([s']))=B([s']).
\]
If \(B([s])=B([s'])=b\), both quotient normal forms lie in \(C_b\). By the unique consistent
extension assumption, \(C_b\) has at most one element, so the quotient normal forms are equal.
\end{proof}

\begin{theorem}[Functional selected-fiber uniqueness]\label{thm:functional-selected-fiber}
Let \(Q_b:=B^{-1}(b)\subseteq\Sigma/\Gamma\) be a same-boundary quotient fiber presented on a
rooted finite overlap graph with spanning tree \(T\), root \(r\), and patch state sets \(X_v\).
Assume the boundary value fixes the root value
\[
u_r=\beta(b),
\]
and every non-root vertex has a deterministic extension map
\[
f_v:X_{p(v)}\times\mathcal B\to X_v .
\]
Define \(u_b\) recursively by
\[
u_v=f_v(u_{p(v)},b).
\]
Evaluate all non-tree overlap equations, sector constraints, and holonomy checks on this
candidate. If they all pass, then \(C_b=\{u_b\}\). If any check fails, then \(C_b=\varnothing\).
\end{theorem}

\begin{proof}
Every consistent state in \(Q_b\) must have root value \(\beta(b)\). Along each tree edge, the
deterministic extension equation forces the child value to be
\(f_v(u_{p(v)},b)\). Induction over tree depth therefore forces every consistent state in the
fiber to equal the single recursively constructed candidate \(u_b\). The remaining constraints are
exactly the non-tree overlaps, sector equations, and holonomy checks. If they pass, \(u_b\) is a
consistent state and no other state can be. If one fails, the only tree-compatible candidate is
inconsistent, so no element of \(C_b\) exists.
\end{proof}

\begin{proposition}[Tree repair realizes the selected extension]\label{prop:tree-selected-descent}
Under the hypotheses of Theorem~\ref{thm:functional-selected-fiber}, choose positive weights
\[
w_v>\sum_{c:p(c)=v}w_c
\]
and set
\[
\Phi_T(x)=\sum_{v\ne r}w_v\,\mathbf 1[x_v\ne f_v(x_{p(v)},b)] .
\]
The repair that sets a non-root vertex to \(f_v(x_{p(v)},b)\) strictly decreases \(\Phi_T\), and
the aggregate transaction for a connected tree-repair conflict component computes the same result
in increasing tree depth.
\end{proposition}

\begin{proof}
Repairing \(v\) removes the \(v\)-term of weight \(w_v\). Only children of \(v\) can become newly
unmatched, contributing at most \(\sum_{c:p(c)=v}w_c\), which is strictly smaller than \(w_v\).
Thus the tree potential decreases. For a connected conflict component, applying parent repairs
before child repairs is forced by the same extension equations, and any different
parenthesization has the same quotient result because each child value is a deterministic function
of the repaired parent and \(b\).
\end{proof}

\begin{corollary}[Nontrivial branch elimination]\label{cor:branch-elimination}
Assume transactional confluence, boundary preservation, repair completeness, and
Theorem~\ref{thm:functional-selected-fiber}. If a selected boundary fiber has \(|Q_b|\ge2\) and
\(C_b=\{u_b\}\), then
\[
\overline{\operatorname{nf}}_\lambda(Q_b)=\{u_b\},
\]
while \(Q_b\setminus C_b\ne\varnothing\). Thus multiple candidate interiors exist, inconsistent
candidates are eliminated by repair, and all surviving candidates share one quotient normal form.
If \(C_b=\varnothing\), the branch is obstructed. If a union solver returns two physically
distinct consistent quotient endpoints for the same selected data, the correct status is
\textsc{ambiguous-union-repair} or an explicit continuation gate, not selection by hash.
\end{corollary}

\begin{remark}[Receipts do not replace the theorem]\label{rem:receipt-boundary}
Simulator receipts named seam descent, atomic commit, distributed local diamond, repair
completeness, same-boundary multistart confluence, and quotient normal-form canonical hash are
public evidence contracts for the finite theorem hypotheses above. They are not proof substitutes.
In particular, a normal-form hash certifies that two emitted quotient normal forms agree after the
declared quotienting; it is not allowed to choose between two physically distinct minimizing
states.
\end{remark}

\begin{corollary}[Gauge-invariant law]\label{cor:gaugeinv}
If $M:\Sigma\to Y$ is gauge-invariant ($M(\gamma\cdot s)=M(s)$ for all $\gamma$), then $M(\operatorname{nf}_\lambda(s))$ depends only on the gauge orbit $[s]$, not the representative.
\end{corollary}

\begin{proof}
Because $M$ is gauge-invariant, it factors through the orbit map: $M=\overline M\circ q$ for some
$\overline M:\Sigma/\Gamma\to Y$. Theorem~\ref{thm:gauge} gives
\[
q\bigl(\operatorname{nf}_\lambda(\gamma\cdot s)\bigr)
=
q\bigl(\operatorname{nf}_\lambda(s)\bigr),
\]
hence
\[
M\bigl(\operatorname{nf}_\lambda(\gamma\cdot s)\bigr)
=
\overline M\!\left(q\bigl(\operatorname{nf}_\lambda(\gamma\cdot s)\bigr)\right)
=
\overline M\!\left(q\bigl(\operatorname{nf}_\lambda(s)\bigr)\right)
=
M\bigl(\operatorname{nf}_\lambda(s)\bigr).
\]
\end{proof}

\begin{definition}[Physical observable algebra on the quantum lift]\label{def:physobs}
Fix a finite patch region or union collar \(R\) on the declared fixed-cutoff quantum lift of
Appendix~\ref{app:consensus-fixed-cutoff-port}. Its \textbf{physical observable algebra}
\(\mathcal A_{\mathrm{phys}}(R)\) is the fixed-point collar algebra under the compact boundary
redundancy action on the ordinary or central-defect branch, or the corresponding quotient-local
algebra on the genuinely noncentral branch. The central sector projectors carried by the collar
decomposition belong to \(Z(\mathcal A_{\mathrm{phys}}(R))\). A \textbf{physical observable} is
any \(X\in \mathcal A_{\mathrm{phys}}(R)\). For a microscopic representative \(s\in\Sigma\), let
\(\omega_R^s\) denote the induced state on \(\mathcal A_{\mathrm{phys}}(R)\).
\end{definition}

\begin{theorem}[Observable-level confluence on the quantum lift]\label{thm:obsconfluence}
Under the hypotheses of Theorems~\ref{thm:confluence} and \ref{thm:gauge}, every initial orbit
\([s]\in\Sigma/\Gamma\) has a unique quotient normal form
\[
\overline{\operatorname{nf}}_\lambda([s])\in q(C).
\]
Fix a declared region \(R\). Let \(t,u\in\Sigma\) be terminal microscopic representatives reached
by representative lifts of maximal repair sequences from initial states in the same orbit
\([s]\). If the induced \(R\)-collar states of \(t\) and \(u\) are representatives of the same
quotient-local glued state in the sense of
Proposition~\ref{prop:assocglue}, then for every physical observable
\(X\in\mathcal A_{\mathrm{phys}}(R)\),
\[
\omega_R^t(X)=\omega_R^u(X).
\]
Hence all physical observables converge to the same overlap-consistent values even when the
microscopic terminal representatives differ by gauge relabelings globally or by sector/higher-gauge
relabelings inside one declared quotient-local glued state.
\end{theorem}

\begin{proof}
Theorems~\ref{thm:confluence} and \ref{thm:gauge} give the unique quotient normal form
\(\overline{\operatorname{nf}}_\lambda([s])=[\operatorname{nf}_\lambda(s)]\in q(C)\). Let
\(t\) and \(u\) be as stated. By Corollary~\ref{cor:obsglue}, representatives of the same
quotient-local glued state induce the same state on the physical observable algebra
\(\mathcal A_{\mathrm{phys}}(R)\), including the central sector projectors carried by that
algebra. Therefore \(\omega_R^t(X)=\omega_R^u(X)\) for every
\(X\in\mathcal A_{\mathrm{phys}}(R)\).
\end{proof}

\begin{remark}[Inputs and boundary for observable-level confluence]\label{rem:obsconfluence}
Theorem~\ref{thm:obsconfluence} does not add a new repair hypothesis beyond the confluence and
fixed-cutoff gluing package. Its inputs are exactly the representative-level confluence theorem,
quotient descent of the repair law, Definition~\ref{def:physobs}, and
Corollary~\ref{cor:obsglue} from the fixed-cutoff union-collar gluing package. The boundary item
is extension of the same statement to broader refinement-stable branches where the declared
union-collar compatibility is only approximate or where the chosen physical observable algebra
itself changes under refinement.
\end{remark}

\begin{corollary}[Inert ancillary refinement does not change physical law]\label{cor:ancilla}
Let \(K=\prod_i K_i\) be a finite ancillary state space and define \(\Sigma^\eta:=\Sigma\times K\). Lift the repair maps by
\[
T_i^{\lambda,\eta}(s,k):=(T_i^\lambda(s),k).
\]
If \(M^\eta:\Sigma^\eta\to Y\) ignores the ancillary factor,
\[
M^\eta(s,k)=M(s),
\]
with \(M\) gauge-invariant on \(\Sigma\), then
\[
M^\eta(\operatorname{nf}_\lambda^\eta(s,k))
=
M(\operatorname{nf}_\lambda(s))
\]
for all \((s,k)\in\Sigma^\eta\).
\end{corollary}

\begin{proof}
Because the ancillary factor is inert,
\[
\operatorname{nf}_\lambda^\eta(s,k)=\bigl(\operatorname{nf}_\lambda(s),k\bigr).
\]
Therefore \(M^\eta(\operatorname{nf}_\lambda^\eta(s,k))=M(\operatorname{nf}_\lambda(s))\), and Corollary~\ref{cor:gaugeinv} supplies gauge-orbit independence.
\end{proof}

Physical uniqueness therefore holds on the quotient by gauge or implementation hiding. The same
statement is unchanged under inert ancillary stabilization.

\begin{definition}[Finite packet closure simplex]\label{def:packet-closure-simplex}
Assume the finite fixed-cutoff consensus branch of Theorems~\ref{thm:confluence} and
\ref{thm:gauge}, so the physical quotient \(Q:=\Sigma/\Gamma\) is finite and carries the
schedule-independent quotient normal-form map
\[
\overline{\operatorname{nf}}_\lambda:Q\to Q.
\]
Let
\[
N_\lambda:=\overline{\operatorname{nf}}_\lambda(Q)=q(C)
\]
be the quotient normal-form set. The \textbf{finite packet simplex} on \(Q\) is
\[
\Delta(Q):=
\left\{
\mu:Q\to\mathbb R_{\ge0}\ \middle|\ \sum_{x\in Q}\mu(x)=1
\right\},
\]
and the \textbf{finite packet closure map} is the pushforward
\[
\mathcal C_\lambda:\Delta(Q)\to\Delta(Q),
\qquad
\mathcal C_\lambda(\mu):=
\bigl(\overline{\operatorname{nf}}_\lambda\bigr)_*\mu,
\]
equivalently
\[
\mathcal C_\lambda(\mu)(y)
=
\sum_{\substack{x\in Q\\ \overline{\operatorname{nf}}_\lambda(x)=y}}\mu(x).
\]
\end{definition}

\begin{theorem}[Finite packet-quotient closure map]\label{thm:finite-packet-closure}
On the finite fixed-cutoff consensus branch of Definition~\ref{def:packet-closure-simplex}, the
map \(\mathcal C_\lambda\) is an affine continuous idempotent self-map of \(\Delta(Q)\). Its
image is exactly the normal-form simplex
\[
\Delta(N_\lambda)
=
\left\{
\mu\in\Delta(Q)\ \middle|\ \operatorname{supp}(\mu)\subseteq N_\lambda
\right\}.
\]
Hence the fixed points of \(\mathcal C_\lambda\) are exactly the packets supported on quotient
normal forms. For every initial quotient state \([s]\in Q\),
\[
\mathcal C_\lambda(\delta_{[s]})
=
\delta_{\overline{\operatorname{nf}}_\lambda([s])}.
\]
\end{theorem}

\begin{proof}
Because \(\overline{\operatorname{nf}}_\lambda:Q\to Q\) is a set map on a finite set, its
pushforward on probability packets is affine and continuous in the finite-dimensional simplex
topology.

By Theorems~\ref{thm:confluence} and \ref{thm:gauge}, the quotient normal form is
schedule-independent and terminal, so
\[
\overline{\operatorname{nf}}_\lambda\circ\overline{\operatorname{nf}}_\lambda
=
\overline{\operatorname{nf}}_\lambda.
\]
Pushforward therefore gives
\[
\mathcal C_\lambda^2
=
\bigl(\overline{\operatorname{nf}}_\lambda\bigr)_*
\bigl(\overline{\operatorname{nf}}_\lambda\bigr)_*
=
\bigl(\overline{\operatorname{nf}}_\lambda\circ\overline{\operatorname{nf}}_\lambda\bigr)_*
=
\mathcal C_\lambda,
\]
so \(\mathcal C_\lambda\) is idempotent.

If \(\nu=\mathcal C_\lambda(\mu)\), then every point in \(\operatorname{supp}(\nu)\) lies in the
image of \(\overline{\operatorname{nf}}_\lambda\), namely \(N_\lambda\). Thus
\(\operatorname{im}(\mathcal C_\lambda)\subseteq \Delta(N_\lambda)\). Conversely, if
\(\nu\in\Delta(N_\lambda)\), then \(\overline{\operatorname{nf}}_\lambda(y)=y\) for every
\(y\in N_\lambda\), so
\[
\mathcal C_\lambda(\nu)=\nu.
\]
Hence \(\Delta(N_\lambda)\subseteq \operatorname{im}(\mathcal C_\lambda)\), proving
\(\operatorname{im}(\mathcal C_\lambda)=\Delta(N_\lambda)\). The same identity shows that
\(\nu\) is a fixed point if and only if it is supported on \(N_\lambda\). The Dirac-mass formula
is the pushforward of a point mass under \(\overline{\operatorname{nf}}_\lambda\).
\end{proof}

\begin{remark}[Boundary of the finite closure theorem]\label{rem:finite-packet-closure}
Theorem~\ref{thm:finite-packet-closure} is an exact finite-branch result. It proves a closure map
only on the finite packet simplex built from one fixed quotient carrier whose repair law is
terminating, confluent, and quotient-descended. It does not prove a habitat-level closure
map for arbitrary OPH state-and-law data, does not identify a nonempty observer-supporting
invariant sector inside the Appendix-B habitat, and does not supply uniqueness or stability
estimates beyond this finite packet branch.
\end{remark}


\section{Refinement-Limit Consensus Classes}
\label{sec:refinement-consensus}

The finite consensus theorem is the operational object used by the continuum branches. To pass
from one finite patch net to a refining family, the paper uses an explicit inverse-limit package
with named compatibility clauses.

\begin{definition}[Separated cofinal refinement consensus system]\label{def:refinement-consensus}
Let \((R,\preceq)\) be a directed refinement set. For each \(r\in R\), let
\[
Q_r:=\Sigma_r/\Gamma_r
\]
be the finite physical quotient state space of a patch presentation, let
\[
n_r:Q_r\to Q_r
\]
be the finite-stage quotient normal-form map supplied by Theorems~\ref{thm:confluence} and
\ref{thm:gauge}, and let
\[
h_r:Q_r\to \mathcal H_r
\]
be the finite-stage holonomy or higher-gauge obstruction map supplied by
Theorems~\ref{thm:holonomy} and \ref{thm:higherdefects}. For \(r\preceq s\), assume restriction
maps
\[
\rho_{sr}:Q_s\to Q_r,
\qquad
\chi_{sr}:\mathcal H_s\to\mathcal H_r,
\]
with \(\rho_{rr}=\mathrm{id}\), \(\chi_{rr}=\mathrm{id}\), and the cocycle identities
\[
\rho_{tr}=\rho_{sr}\circ\rho_{ts},
\qquad
\chi_{tr}=\chi_{sr}\circ\chi_{ts}
\qquad
(r\preceq s\preceq t).
\]
The system is a separated cofinal refinement consensus system when:
\begin{enumerate}
\item \textbf{Normal-form naturality:}
\[
\rho_{sr}\circ n_s=n_r\circ\rho_{sr}
\qquad
(r\preceq s).
\]
\item \textbf{Holonomy naturality:}
\[
\chi_{sr}\circ h_s=h_r\circ\rho_{sr}
\qquad
(r\preceq s).
\]
\item \textbf{Visible separation:} two compatible families in \(\varprojlim Q_r\), or in
\(\varprojlim \mathcal H_r\), are equal whenever their projections agree on a cofinal subset of
stages.
\end{enumerate}
\end{definition}

\begin{theorem}[Refinement-limit consensus and holonomy classes]\label{thm:refinement-consensus}
For any separated cofinal refinement consensus system, the formulas
\[
n_\infty\bigl((x_r)_r\bigr):=(n_r(x_r))_r,
\qquad
h_\infty\bigl((x_r)_r\bigr):=(h_r(x_r))_r
\]
define maps
\[
n_\infty:\varprojlim Q_r\to\varprojlim Q_r,
\qquad
h_\infty:\varprojlim Q_r\to\varprojlim\mathcal H_r.
\]
The class \(n_\infty(x)\) is the unique schedule-independent refinement-limit normal form of
\(x\). The class \(h_\infty(x)\) is the refinement-limit holonomy obstruction. The pair
\((n_\infty(x),h_\infty(x))\) is the refinement-limit consensus class. In the inverse-limit
topology, finite-stage normal forms and holonomy classes converge to these two classes: for every
stage \(r\), all refinements \(s\succeq r\) restrict to the fixed values
\[
\rho_{sr}(n_s(x_s))=n_r(x_r),
\qquad
\chi_{sr}(h_s(x_s))=h_r(x_r).
\]
Furthermore, \(h_\infty(x)=0\) if and only if every finite projection has zero finite-stage
obstruction. If \(h_\infty(x)\ne0\), visible separation gives a finite stage that witnesses the
nonzero obstruction. If two refinement-limit candidates have the same normal-form and holonomy
projections on a cofinal tail, then they determine the same refinement-limit consensus class.
\end{theorem}

\begin{proof}
Let \(x=(x_r)_r\in\varprojlim Q_r\), so \(\rho_{sr}(x_s)=x_r\) for \(r\preceq s\). Normal-form
naturality gives
\[
\rho_{sr}(n_s(x_s))
=
n_r(\rho_{sr}(x_s))
=
n_r(x_r),
\]
so \((n_r(x_r))_r\) is a compatible family and \(n_\infty\) is well defined. The same argument
with holonomy naturality gives
\[
\chi_{sr}(h_s(x_s))
=
h_r(\rho_{sr}(x_s))
=
h_r(x_r),
\]
so \(h_\infty\) is well defined.

Each finite \(n_r(x_r)\) is independent of update schedule by Theorems~\ref{thm:confluence} and
\ref{thm:gauge}. Therefore every finite projection of \(n_\infty(x)\) is schedule-independent,
and visible separation makes the inverse-limit class unique. The displayed restriction identities
are exactly the cylinder convergence statement in the inverse-limit topology.

The zero-obstruction statement follows from the definition of the inverse-limit zero class. The
identity \(h_\infty(x)=0\) holds exactly when every projection \(h_r(x_r)\) is zero. If
\(h_\infty(x)\ne0\), at least one finite projection is nonzero, and any cofinal tail containing a
refinement of that stage carries the same nonzero projected obstruction by holonomy naturality.
The last claim is visible separation applied to the compatible normal-form and holonomy families.
\end{proof}

\begin{remark}[Scope of the refinement theorem]\label{rem:refinement-consensus}
Theorem~\ref{thm:refinement-consensus} is a controlled continuum bridge. It proves persistence,
exhaustion, and collapse of the consensus data once the OPH refinement system supplies the
restriction maps and the two naturality clauses in Definition~\ref{def:refinement-consensus}. It
leaves uniform complexity bounds, automatic fair-block noisy-consensus certificates, and
automatic normal-form naturality for arbitrary changing repair laws as separate branch conditions.
\end{remark}

\begin{definition}[Repair morphism]\label{def:repair-morphism}
Let \((Q_s,\to_s,C_s,n_s)\) and \((Q_r,\to_r,C_r,n_r)\) be two finite quotient repair systems
covered by Theorems~\ref{thm:confluence} and \ref{thm:gauge}. A map
\[
\rho_{sr}:Q_s\to Q_r
\]
is a \textbf{repair morphism} when:
\begin{enumerate}[label=(\roman*),leftmargin=2em]
\item every fine repair step \(x\to_s y\) maps to a coarse repair path or a stutter,
\[
\rho_{sr}(x)\to_r^*\rho_{sr}(y);
\]
\item fine consistency maps into coarse consistency,
\[
\rho_{sr}(C_s)\subseteq C_r.
\]
\end{enumerate}
For a concrete local repair law this is checked by giving, for every fine repair generator, a
coarse witness word in the accepted coarse repair generators or the empty word, and by verifying
that the restriction maps carry every fine consistency equation to a coarse consistency equation.
\end{definition}

\begin{theorem}[Normal-form naturality from repair morphisms]\label{thm:repair-morphism-naturality}
If \(\rho_{sr}:Q_s\to Q_r\) is a repair morphism, then
\[
\rho_{sr}\circ n_s=n_r\circ\rho_{sr}.
\]
\end{theorem}

\begin{proof}
Fix \(x\in Q_s\). Since \(x\to_s^* n_s(x)\), repeated use of the step-simulation clause gives
\[
\rho_{sr}(x)\to_r^*\rho_{sr}(n_s(x)).
\]
The consistency-preservation clause gives \(\rho_{sr}(n_s(x))\in C_r\), because
\(n_s(x)\in C_s\). Thus \(\rho_{sr}(n_s(x))\) is a coarse normal form reachable from
\(\rho_{sr}(x)\). The coarse system has a unique normal form reachable from each initial quotient
state, so \(\rho_{sr}(n_s(x))=n_r(\rho_{sr}(x))\).
\end{proof}

\begin{definition}[Holonomy cochain morphism]\label{def:holonomy-cochain-morphism}
For obstruction maps \(h_s:Q_s\to\mathcal H_s\) and \(h_r:Q_r\to\mathcal H_r\), a map
\(\chi_{sr}:\mathcal H_s\to\mathcal H_r\) is a \textbf{holonomy cochain morphism} when the
edge, cycle, or crossed-module cochains used to compute \(h_s\) restrict to the cochains used to
compute \(h_r\), modulo the declared quotient identifications. Equivalently,
\[
\chi_{sr}\circ h_s=h_r\circ\rho_{sr}.
\]
\end{definition}

\begin{remark}[How exact refinement naturality is discharged]\label{rem:exact-naturality-discharge}
Definition~\ref{def:refinement-consensus} may be used as an interface contract, but on an
implemented exact branch the normal-form part is proved by Theorem~\ref{thm:repair-morphism-naturality}
from explicit coarse witness words for fine repairs. The holonomy part is proved by the cochain
morphism check of Definition~\ref{def:holonomy-cochain-morphism}. Approximate RG branches use the
controlled defects of Definition~\ref{def:controlled-coarse-consensus} instead.
\end{remark}

\begin{definition}[Controlled coarse-graining / reconciliation square]\label{def:controlled-coarse-consensus}
Fix refinement stages \(r\preceq s\). Let \(Q_r,Q_s\) be the physical quotient state spaces,
\(n_r,n_s\) their finite-stage normal-form maps, and \(h_r,h_s\) their holonomy or
higher-gauge obstruction maps. Let
\[
\rho_{sr}:Q_s\to Q_r,
\qquad
\chi_{sr}:\mathcal H_s\to\mathcal H_r
\]
be the coarse-graining maps on quotient states and obstruction data. Equip \(Q_r\) and
\(\mathcal H_r\) with pseudometrics \(d^Q_r\) and \(d^{\mathcal H}_r\) that define the
macroscopic readout scale at stage \(r\).

The square is \((\varepsilon^n_{sr},\varepsilon^h_{sr})\)-controlled when, for every
\(x\in Q_s\),
\[
d^Q_r\!\left(\rho_{sr}(n_s(x)),\,n_r(\rho_{sr}(x))\right)
\le \varepsilon^n_{sr},
\]
and
\[
d^{\mathcal H}_r\!\left(\chi_{sr}(h_s(x)),\,h_r(\rho_{sr}(x))\right)
\le \varepsilon^h_{sr}.
\]
The errors are \textbf{cofinally vanishing} when, for every fixed macroscopic stage \(r\) and
every \(\delta>0\), there is a refinement stage \(s_0\succeq r\) such that
\(\varepsilon^n_{sr},\varepsilon^h_{sr}<\delta\) for all \(s\succeq s_0\).
\end{definition}

\begin{theorem}[Coarse-graining commutes with reconciliation up to controlled error]\label{thm:coarse-grain-consensus}
Suppose the refinement square of Definition~\ref{def:controlled-coarse-consensus} is
\((\varepsilon^n_{sr},\varepsilon^h_{sr})\)-controlled. Define the coarse-stage consensus readout
\[
\mathcal C_r(y):=(n_r(y),h_r(y))
\]
and the fine-then-coarse readout
\[
\mathcal C_{s\to r}^{\mathrm{fine}}(x):=(\rho_{sr}(n_s(x)),\chi_{sr}(h_s(x))).
\]
Then, for every fine state \(x\in Q_s\), the product readout distance between
\(\mathcal C_{s\to r}^{\mathrm{fine}}(x)\) and \(\mathcal C_r(\rho_{sr}(x))\) is bounded by
\[
\max\{\varepsilon^n_{sr},\varepsilon^h_{sr}\}.
\]
Thus reconciling at the fine stage and then coarse-graining gives the same macroscopic law data
as coarse-graining first and reconciling at the coarse stage, up to the declared control errors.
If the exact naturality clauses of Definition~\ref{def:refinement-consensus} hold, then
\(\varepsilon^n_{sr}=\varepsilon^h_{sr}=0\). If the errors are cofinally vanishing, the two
procedures determine the same cylinder values in the inverse-limit topology at every fixed
macroscopic stage.
\end{theorem}

\begin{proof}
The normal-form component of the product readout is bounded by the first inequality in
Definition~\ref{def:controlled-coarse-consensus}, and the obstruction component is bounded by the
second. Taking the maximum gives the stated product bound. Exact naturality is precisely the
special case
\[
\rho_{sr}\circ n_s=n_r\circ\rho_{sr},
\qquad
\chi_{sr}\circ h_s=h_r\circ\rho_{sr},
\]
so both errors vanish there. Cofinal vanishing means that, for any fixed coarse cylinder and any
desired readout tolerance, all sufficiently fine presentations give the same cylinder value
within that tolerance, which is convergence in the inverse-limit topology.
\end{proof}

\begin{remark}[What remains to verify on a concrete RG branch]\label{rem:coarse-grain-scope}
Theorem~\ref{thm:coarse-grain-consensus} is not a claim that arbitrary renormalization maps
commute with arbitrary repair laws. It identifies the exact branch burden: derive or estimate the
two square defects \(\varepsilon^n_{sr}\) and \(\varepsilon^h_{sr}\) from the chosen
coarse-graining channel, recovery map, decoder, and obstruction readout. The exact refinement
theorem is the zero-defect case; approximate RG matching is theorem-grade only when these defects
are explicitly controlled.
\end{remark}


\section{Record Algebras and the Operator Observation Layer}

\paragraph{Inputs used here.}
From the fixed-cutoff collar package we use only the observer-accessible finite-dimensional algebra
on one completed compare/write/verify slice, the declared pointer and overlap-sector projectors
read on that same slice, and the trace-distance control on restored accessible states when
restoration is invoked. No continuum lift and no broader observer-metaphysical assumption is used
here.

\begin{definition}[Exact record algebras and approximate record presentations]\label{def:lattice}
Fix one completed observer-accessible slice at cycle \(t\), and let
\(\mathcal A_t^{\mathrm{acc}}\) be the corresponding finite-dimensional accessible algebra. An
\emph{exact record presentation} is a family of orthogonal projectors
\(\{\widehat P_a(t)\}_a\subset Z(\mathcal A_t^{\mathrm{acc}})\) whose generated algebra
\[
\mathcal Z_{\mathrm{rec}}(t)
:=
\mathrm{Alg}\bigl(\{\widehat P_a(t)\}_a\bigr)
\]
is finite and commutative.

An \emph{approximate record presentation} on the same declared readout slots is a family of
projectors \(\{P_a(t)\}_a\subset \mathcal A_t^{\mathrm{acc}}\) together with an exact record
presentation \(\{\widehat P_a(t)\}_a\) such that
\[
\delta_{\mathrm{rec}}(t)
:=
\max_a \|P_a(t)-\widehat P_a(t)\|
\]
is finite.
\end{definition}

\begin{theorem}[Record algebra, Born-L\"uders update, and quantitative stability]\label{thm:crdt}
Fix one completed observer-accessible slice at cycle \(t\).
\begin{enumerate}
\item The exact record projectors generate a finite commutative central algebra
\(\mathcal Z_{\mathrm{rec}}(t)\subset Z(\mathcal A_t^{\mathrm{acc}})\).
\item For every event \(E\) in the finite event algebra generated by \(\mathcal Z_{\mathrm{rec}}(t)\),
\[
\mathbb P_t(E)=\operatorname{Tr}\!\bigl(\rho_t \widehat P_E(t)\bigr),
\qquad
\rho_t\!\mid_E
=
\frac{\widehat P_E(t)\rho_t \widehat P_E(t)}
{\operatorname{Tr}(\rho_t \widehat P_E(t))}.
\]
\item If no accepted repair between \(t\) and \(t+1\) touches the support of \(\widehat P_E(t)\),
then the next read of \(E\) has probability \(1\).
\item If \(\{P_a(t)\}_a\) is an approximate record presentation with modulus
\(\delta_{\mathrm{rec}}(t)\), then
\[
\|[P_a(t),P_b(t)]\|\le 4\,\delta_{\mathrm{rec}}(t)
\]
for all declared record projectors \(P_a(t),P_b(t)\). Moreover, for every declared elementary
record event \(a\) and every restored accessible state \(\widetilde\rho_t\) satisfying
\[
\|\widetilde\rho_t-\rho_t\|_1\le \varepsilon,
\]
one has
\[
\Bigl|
\operatorname{Tr}\!\bigl(\widetilde\rho_t P_a(t)\bigr)
-
\operatorname{Tr}\!\bigl(\rho_t \widehat P_a(t)\bigr)
\Bigr|
\le
\varepsilon+\delta_{\mathrm{rec}}(t).
\]
\end{enumerate}
\end{theorem}

\begin{proof}
Because the exact record projectors lie in the center of \(\mathcal A_t^{\mathrm{acc}}\), they
generate a finite commutative central algebra. Finite-dimensional projective measurement on that
commuting algebra gives the Born trace and the L\"uders conditioned state, proving (1) and (2).

If no accepted repair touches the support of \(\widehat P_E(t)\), then the next completed read is
performed on the same central projector surface. After conditioning on \(E\), the state lies in
the range of \(\widehat P_E(t)\), so the next read of the same event has probability \(1\). This
gives (3).

For (4), write
\[
[P_a(t),P_b(t)]
=
[P_a(t)-\widehat P_a(t),P_b(t)]
+
[\widehat P_a(t),P_b(t)-\widehat P_b(t)],
\]
because \([\widehat P_a(t),\widehat P_b(t)]=0\). Since every projector has operator norm at most
\(1\),
\[
\|[P_a(t),P_b(t)]\|
\le
2\|P_a(t)-\widehat P_a(t)\|
+
2\|P_b(t)-\widehat P_b(t)\|
\le
4\,\delta_{\mathrm{rec}}(t).
\]
For the probability bound,
\[
\Bigl|
\operatorname{Tr}\!\bigl(\widetilde\rho_t P_a(t)\bigr)
-
\operatorname{Tr}\!\bigl(\rho_t \widehat P_a(t)\bigr)
\Bigr|
\le
\Bigl|\operatorname{Tr}\!\bigl((\widetilde\rho_t-\rho_t)P_a(t)\bigr)\Bigr|
+
\Bigl|\operatorname{Tr}\!\bigl(\rho_t(P_a(t)-\widehat P_a(t))\bigr)\Bigr|.
\]
The first term is bounded by \(\|\widetilde\rho_t-\rho_t\|_1\|P_a(t)\|\le \varepsilon\), and the
second by \(\|\rho_t\|_1\|P_a(t)-\widehat P_a(t)\|\le \delta_{\mathrm{rec}}(t)\). Hence the total
error is at most \(\varepsilon+\delta_{\mathrm{rec}}(t)\).
\end{proof}

\paragraph{Merge boundary.}
What is closed here is the fixed-cutoff operator-algebraic observation surface: a central record
algebra on the exact readout slice, Born/L\"uders measurement on its event projectors, exact
repeated-read stability under the untouched-support hypothesis, and explicit
\((\varepsilon,\delta_{\mathrm{rec}})\) control when practical readout projectors are only close
to that central surface. The broader export problem asks whether richer branches keep physically
relevant pointer surfaces close to one such central record algebra and whether the same control
survives refinement and continuation.

\section{OPH Simulation Criterion and Fixed-Point Firewall}
\label{sec:oph-simulation}

A fixed-point equation is necessary for a selected finite OPH packet, but it is not the
definition of an OPH simulation. The additional clauses below are the nontrivial content.

\begin{definition}[Nontrivial OPH simulation certificate]\label{def:oph-simulation}
Work on the finite fixed-cutoff branch of Definition~\ref{def:packet-closure-simplex}. Write
\[
 Q:=\Sigma/\Gamma,
 \qquad
 n:=\overline{\operatorname{nf}}_\lambda:Q\to Q,
 \qquad
 N_\lambda:=q(C).
\]
Let \(B:Q\to\mathcal B\) be a boundary/sector map preserved by accepted quotient repairs,
fix \(b\in\mathcal B\), and set
\[
 Q_b:=B^{-1}(b),
 \qquad
 C_b:=N_\lambda\cap Q_b.
\]
For \(u\in C_b\), let \(\mathfrak U_u:=\delta_u\in\Delta(Q)\). The selected pair
\((b,u)\), equivalently the packet \(\mathfrak U_u\), has an \textbf{OPH simulation
certificate} when all of the following hold.
\begin{enumerate}[label=(S\arabic*),leftmargin=3.0em]
\item \textbf{Endogenous update.} The local maps are the recovery-derived collar maps of
Definition~\ref{def:repair}, descend to the physical quotient, and use only the declared
patch, overlap, recovery, decoder, and acceptance data. No undeclared oracle variable is an
argument of an accepted physical update.

\item \textbf{Observer-readable records.} A declared observer-accessible region carries an
exact record presentation \(\{\widehat P_a\}_a\) as in Definition~\ref{def:lattice}, with at
least two nonzero orthogonal event projectors. Event probabilities, conditioned readouts, and
any checkpoint/order relation interpreted as history are collected in a physical map
\[
 \operatorname{Read}:Q\to\mathcal Y_{\rm rec}
\]
whose terminal value depends only on the quotient normal form.

\item \textbf{Overlap repair.} Every accepted nontrivial move strictly lowers the declared
touched-overlap score, normal forms are exactly the globally consistent states, and the
quotient normal form \(n\) is terminating and schedule-independent.

\item \textbf{Nontrivial branch elimination.} The selected fiber has one consistent
quotient state but more than one candidate:
\[
 C_b=\{u\},
 \qquad
 Q_b\setminus C_b\ne\varnothing,
 \qquad
 n(Q_b)=\{u\}.
\]
Thus at least one inconsistent candidate is genuinely removed rather than merely renamed a
fixed point. A candidate with a nonzero declared holonomy or higher-gauge obstruction is
rejected as obstructed and is not counted as an alternative selected world.

\item \textbf{Implementation and clock closure.} The physical update and record readout
have types
\[
 n:Q\to Q,
 \qquad
 \operatorname{Read}:Q\to\mathcal Y_{\rm rec},
\]
with no indispensable machine-state or external-time argument. More precisely, for any
auxiliary implementation and clock sets \(\mathcal E_{\rm ext}\) and \(\Theta_{\rm ext}\),
let
\[
 p:Q\times\mathcal E_{\rm ext}\times\Theta_{\rm ext}\to Q
\]
be the physical projection. An admissible lifted update \(\widetilde n\) and readout
\(\widetilde{\operatorname{Read}}\) must satisfy
\[
 p\circ\widetilde n=n\circ p,
 \qquad
 \widetilde{\operatorname{Read}}=\operatorname{Read}\circ p.
\]
The asynchronous schedule and repair-step counter are proof data, not physical clock
coordinates. Any history claimed by the branch is read from records in the terminal quotient
state rather than supplied by an external clock.
\end{enumerate}
The identities
\[
 n(u)=u,
 \qquad
 \mathcal C_\lambda(\mathfrak U_u)=\mathfrak U_u
\]
are therefore consequences of the certificate, not its definition.
\end{definition}

\begin{theorem}[Selected OPH packet and fixed-point firewall]\label{thm:oph-simulation}
Assume Definitions~\ref{def:repair}, \ref{def:localfit}, and \ref{def:assoc},
Assumption~\ref{ass:complete}, and the quotient and record hypotheses of
Theorems~\ref{thm:gauge}, \ref{thm:obsconfluence}, and \ref{thm:crdt}. Fix a preserved
boundary/sector value \(b_{\rm OPH}\) for which
\[
 C_{b_{\rm OPH}}=\{u_{\rm OPH}\},
 \qquad
 Q_{b_{\rm OPH}}\setminus C_{b_{\rm OPH}}\ne\varnothing,
\]
and assume the selected observer surface has a nontrivial exact record presentation whose
checkpoint/order data, when interpreted as history, are quotient-observable. Then the selected
finite OPH packet
\[
 \mathfrak U_{\rm OPH}:=\delta_{u_{\rm OPH}}
\]
has an OPH simulation certificate in the sense of Definition~\ref{def:oph-simulation}. In
particular,
\[
 n(u_{\rm OPH})=u_{\rm OPH},
 \qquad
 \mathcal C_\lambda(\mathfrak U_{\rm OPH})=\mathfrak U_{\rm OPH}.
\]

The converse is false: a fixed point, equilibrium, stationary point, or variational solution
does not by itself imply an OPH simulation certificate. Indeed, for any set \(X\) with
\(|X|>1\), the identity map \(F=\mathrm{id}_X\) satisfies \(F(x)=x\) for every \(x\in X\),
and every \(x\) is a global minimizer and stationary point of the constant functional
\(V\equiv0\). Nevertheless, \(F^{-1}(x)=\{x\}\), so no proper candidate basin collapses to
\(x\); clause \emph{(S4)} fails. The fixed-point or variational equation alone also supplies
none of clauses \emph{(S1)}--\emph{(S3)} or \emph{(S5)}.
\end{theorem}

\begin{proof}
Clause \emph{(S1)} is Definition~\ref{def:repair} together with quotient descent from
Theorem~\ref{thm:gauge}. Clause \emph{(S2)} follows from Theorem~\ref{thm:crdt}; terminal
representative independence is Theorem~\ref{thm:obsconfluence}. Definition~\ref{def:localfit},
Proposition~\ref{prop:lyapunov-termination}, Proposition~\ref{prop:local-confluence},
Assumption~\ref{ass:complete}, and Theorem~\ref{thm:confluence} give clause \emph{(S3)}.

For \(x\in Q_{b_{\rm OPH}}\), boundary preservation gives
\(B(n(x))=b_{\rm OPH}\), while repair completeness gives \(n(x)\in N_\lambda\). Hence
\(n(x)\in C_{b_{\rm OPH}}=\{u_{\rm OPH}\}\), proving
\(n(Q_{b_{\rm OPH}})=\{u_{\rm OPH}\}\). The assumed nonempty difference
\(Q_{b_{\rm OPH}}\setminus C_{b_{\rm OPH}}\) makes this elimination nontrivial. The
holonomy and higher-gauge rejection statement is Theorems~\ref{thm:holonomy} and
\ref{thm:higherdefects}. This proves clause \emph{(S4)}. Its proper-basin requirement is
nonvacuous: on the verified rooted-tree packet domain of Theorem~\ref{thm:tree-packet-domain},
fixing the root packet gives a singleton consistent fiber, while changing any non-root packet
supplies an inconsistent candidate in that fiber.

The quotient normal-form and physical-record maps are defined on the declared physical state
alone. Schedule independence removes the scheduler from the output, and
Corollary~\ref{cor:ancilla} proves invariance under inert carrier enlargement. The displayed
projection identities are the admissibility condition for any more general carrier or clock
realization. Thus an auxiliary implementation state or iteration counter cannot change the
physical result, proving clause \emph{(S5)}. Finally,
Theorem~\ref{thm:finite-packet-closure} gives
\[
 \mathcal C_\lambda(\delta_{u_{\rm OPH}})
 =\delta_{n(u_{\rm OPH})}
 =\delta_{u_{\rm OPH}},
\]
because \(u_{\rm OPH}\in N_\lambda\). The identity-map and constant-functional example
proves the final non-implication.
\end{proof}

\begin{remark}[Structural falsification hooks]\label{rem:oph-simulation-falsifiers}
An asserted OPH simulation certificate fails if any one of the following occurs: an accepted
update depends on undeclared oracle, carrier, or external-clock data; the observer record
algebra is trivial, unstable, or not quotient-readable; an accepted move fails strict mismatch
descent; terminal states are not exactly the overlap-consistent states; two admissible schedules
or two candidates in the selected boundary fiber produce different physical normal forms; the
selected fiber has no inconsistent candidate to eliminate; a claimed global branch has nonzero
holonomy or higher-gauge obstruction; or a carrier/clock lift changes the physical update or
readout after projection. These are theorem-level failure conditions, not semantic disagreements
about the word ``simulation.''
\end{remark}

\begin{remark}[Scope]\label{rem:oph-simulation-scope}
Theorem~\ref{thm:oph-simulation} certifies the named finite fixed-cutoff branch and selected
boundary/sector fiber. It does not promote the finite packet closure to an arbitrary habitat-level
state-and-law space, and it does not remove the branch and record hypotheses stated in the
theorem.
\end{remark}

\section{Distributed Presentations of One Finite Universe}
\label{sec:distributed-one-universe}

The implementation closure clause in Definition~\ref{def:oph-simulation} has a distributed
form. A worker partition is allowed to accelerate or package the computation, but it is not a
new physical ingredient. The physical question is whether the worker run presents the same
finite quotient repair system as the monolithic carrier.

\begin{definition}[Finite global OPH carrier]\label{def:finite-global-carrier}
At refinement stage \(r\), a finite global OPH carrier consists of a finite patch graph
\[
G_r=(V_r,E_r),
\]
finite patch state sets \(S_i\), interface alphabets \(I_e\), endpoint maps
\(\pi_{i,e}:S_i\to I_e\), an implementation-hiding group \(\Gamma_r\), a boundary/sector map
\(B_r\), a mismatch potential \(\Phi_r\), accepted quotient repairs \(\to_r\), and an
observer-readable map
\[
\operatorname{Read}_r:Q_r\to\mathcal Y_r,
\qquad
Q_r:=\left(\prod_{i\in V_r}S_i\right)/\Gamma_r.
\]
Let \(C_r\subseteq Q_r\) be the quotient consistency set. On the finite branch covered by
Theorems~\ref{thm:confluence} and \ref{thm:gauge}, the accepted quotient repair relation has a
unique normal-form map
\[
n_r:Q_r\to C_r.
\]
A one-universe input is the pair \(\mathbf U_r=(\mathfrak U_r,q_0)\), where
\(\mathfrak U_r\) is the carrier data above and \(q_0\in Q_r\). A worker partition is not part
of \(\mathfrak U_r\).
\end{definition}

\begin{definition}[Worker presentation and physical projection]\label{def:worker-presentation}
Let \(\kappa:V_r\to W\) assign every patch to one authoritative worker. The cut set is
\[
E_\kappa^{\mathrm{cut}}
:=
\bigl\{\{i,j\}\in E_r:\kappa(i)\ne\kappa(j)\bigr\}.
\]
A distributed presentation for \(\kappa\) has worker-owned states, ghost or halo copies,
message queues, transaction records, retries, worker identifiers, checkpoints, event logs, and
scheduler state. Let \(\widetilde Q_{r,\kappa}\) be its quotient by purely local worker
bookkeeping that leaves authoritative physical states unchanged. The physical projection
\[
p_{r,\kappa}:\widetilde Q_{r,\kappa}\to Q_r
\]
keeps only the authoritative patch states, quotients by \(\Gamma_r\), and discards queues, retry
metadata, worker labels, and external clocks. Let
\[
D_r:=\bigsqcup_{\kappa}\{\kappa\}\times\widetilde Q_{r,\kappa},
\qquad
P_r(\kappa,\tilde q):=p_{r,\kappa}(\tilde q).
\]
A distributed readout is physical when it factors as
\[
\widetilde{\operatorname{Read}}_r=\operatorname{Read}_r\circ P_r.
\]
\end{definition}

\begin{definition}[Admissible distributed event]\label{def:admissible-distributed-event}
For \(d,d'\in D_r\), a distributed event \(d\Rightarrow d'\) is \textbf{admissible} when one of
the following holds.
\begin{enumerate}[label=(\alph*),leftmargin=2em]
\item \textbf{Linearizable physical commit.} There is a nonempty monolithic accepted repair path
\[
P_r(d)\to_r^*P_r(d')
\]
whose repair word is recorded as the event's linearization witness.
\item \textbf{Physical stutter.} \(P_r(d')=P_r(d)\). Message delivery, halo refresh, prepare,
abort, checkpoint, worker start or stop, idempotent replay, and repartition metadata are allowed
only in this class unless they also carry a physical commit witness.
\item \textbf{Certified rollback.} There is an earlier committed state \(d_j\) in the same run
history such that \(P_r(d')=P_r(d_j)\), and the rollback certificate names that committed
projection root. Transparent restart from a committed frontier is a stutter.
\end{enumerate}
\end{definition}

\begin{proposition}[Normal form is constant on an accepted reachability cone]\label{lem:nf-reachability-cone}
If \(q\to_r^*q'\), then
\[
n_r(q')=n_r(q).
\]
\end{proposition}

\begin{proof}
The state \(q'\) is reachable from \(q\). The state \(n_r(q')\) is a normal form reachable from
\(q'\), hence also reachable from \(q\). Theorem~\ref{thm:confluence} and
Theorem~\ref{thm:gauge} give the unique quotient normal form reachable from \(q\), so
\(n_r(q')=n_r(q)\).
\end{proof}

\begin{theorem}[Distributed realization of one finite OPH universe]\label{thm:distributed-one-universe}
Fix a finite global carrier \(\mathbf U_r=(\mathfrak U_r,q_0)\) satisfying
Theorems~\ref{thm:confluence} and \ref{thm:gauge}. Let
\[
d_0\Rightarrow d_1\Rightarrow\cdots\Rightarrow d_m
\]
be a finite distributed execution in \(D_r\) with \(P_r(d_0)=q_0\), and assume every event is
admissible in the sense of Definition~\ref{def:admissible-distributed-event}. Then, for every
\(t\),
\[
q_0\to_r^*P_r(d_t),
\qquad
n_r(P_r(d_t))=n_r(q_0).
\]
If the final projection is a monolithic normal form, then
\[
P_r(d_m)=n_r(q_0).
\]
For every physical observable or observer readout \(M:Q_r\to Y\),
\[
M(P_r(d_m))=M(n_r(q_0))
\]
whenever \(P_r(d_m)\) is normal. In particular,
\(\widetilde{\operatorname{Read}}_r(d_m)=\operatorname{Read}_r(n_r(q_0))\) for every
distributed readout that factors through \(P_r\).
\end{theorem}

\begin{proof}
Induct on \(t\). The claim is true at \(t=0\). Suppose it holds at \(t\). For a linearizable
physical commit, admissibility gives
\[
P_r(d_t)\to_r^*P_r(d_{t+1}),
\]
so \(q_0\to_r^*P_r(d_{t+1})\), and Lemma~\ref{lem:nf-reachability-cone} gives
\[
n_r(P_r(d_{t+1}))=n_r(P_r(d_t))=n_r(q_0).
\]
For a physical stutter, \(P_r(d_{t+1})=P_r(d_t)\), so both statements are unchanged. For a
certified rollback, \(P_r(d_{t+1})\) is the projection of an earlier committed state; the
induction hypothesis already proved reachability from \(q_0\) and equality of normal forms for
that projection. This proves the two displayed identities for every \(t\).

If \(P_r(d_m)\) is a monolithic normal form, then it is the unique quotient normal form reachable
from \(q_0\), hence \(P_r(d_m)=n_r(q_0)\). The observable and readout statements follow by applying
\(M\) or \(\operatorname{Read}_r\) to this equality and using
\(\widetilde{\operatorname{Read}}_r=\operatorname{Read}_r\circ P_r\).
\end{proof}

\begin{corollary}[Partition, schedule, and restart invariance]\label{cor:distributed-invariance}
Any two admissible distributed executions of the same finite carrier \(\mathbf U_r\), with the
same initial quotient state \(q_0\), have the same terminal quotient observables and
observer-readable outputs once their final projections are monolithic normal forms. The statement
is independent of partition \(\kappa\), worker count, scheduler choices, restart history, and
repartition metadata.
\end{corollary}

\begin{proof}
Apply Theorem~\ref{thm:distributed-one-universe} to each execution. Both final projections equal
\(n_r(q_0)\), so every quotient observable and every readout factoring through the projection has
the same value.
\end{proof}

\begin{proposition}[Restart stabilization separates safety from liveness]\label{prop:restart-stabilization}
Assume \(Q_r\) is finite, every physical commit strictly descends the accepted potential until
normality, every rollback returns to an earlier committed projection, only finitely many
progress-erasing rollbacks occur, and after the last such rollback the scheduler is fair enough
to commit an enabled repair whenever the projected state is not normal. Then the distributed run
reaches a monolithic normal form after finitely many physical commits.
\end{proposition}

\begin{proof}
Safety is Theorem~\ref{thm:distributed-one-universe}. For liveness, ignore stutters and the
finitely many progress-erasing rollbacks before the last one. After that time, each nonnormal
projected state eventually takes a strict descending accepted repair. The set of possible
potential values below the post-rollback value is finite, so strict descent can occur only
finitely many times before a normal projected state is reached.
\end{proof}

\begin{theorem}[Distributed refinement cube]\label{thm:distributed-refinement-cube}
Let \(r\preceq s\). Suppose \(\rho_{sr}:Q_s\to Q_r\) is a repair morphism, the corresponding
holonomy maps form a cochain morphism, and distributed presentations at both stages are exact in
the sense of Theorem~\ref{thm:distributed-one-universe}. If the lifted restriction
\(\widetilde\rho_{sr}:D_s\to D_r\) commutes with physical projection,
\[
P_r\circ\widetilde\rho_{sr}=\rho_{sr}\circ P_s,
\]
then distributed execution and coarse restriction commute on normal forms and readouts:
\[
P_r\!\left(\widetilde\rho_{sr}(d_m^s)\right)
=
n_r\!\left(\rho_{sr}(q_0^s)\right)
=
\rho_{sr}\!\left(n_s(q_0^s)\right)
\]
whenever the fine final projection is normal. The same statement holds for holonomy readouts via
\(\chi_{sr}\).
\end{theorem}

\begin{proof}
The fine distributed theorem gives \(P_s(d_m^s)=n_s(q_0^s)\). Projection compatibility gives
\[
P_r(\widetilde\rho_{sr}(d_m^s))=\rho_{sr}(n_s(q_0^s)).
\]
Theorem~\ref{thm:repair-morphism-naturality} identifies this value with
\(n_r(\rho_{sr}(q_0^s))\). The holonomy statement is exactly the cochain-morphism identity of
Definition~\ref{def:holonomy-cochain-morphism}.
\end{proof}

\paragraph{Public certificate contract.}
A run pack that claims Theorem~\ref{thm:distributed-one-universe} must emit enough evidence for
a verifier to reconstruct the premises without trusting worker narration. The required fields are:
a global carrier manifest and hash, the monolithic graph \(G_r\), the global initial quotient
state \(q_0\), a partition map \(\kappa\), cut-interface records \(E_\kappa^{\mathrm{cut}}\)
with restriction maps, a global observer registry, code/config/run hashes, a committed event log
with a linearization witness for each physical commit, stutter records for noncommitting worker
events, rollback records naming earlier committed projection roots, repartition records preserving
the physical projection root, a final monolithic normal-form certificate, and a final readout
recomputed from the projected state. Missing, stale, shard-local, or synthetic replacements for
these fields fail closed. The exact branch uses this contract; noisy branches additionally need
the fair-block constants of Theorem~\ref{thm:noisy-fair-block-consensus}.

\section{Quotient Chart Transport and Neutral Geometry}
\label{sec:quotient-neutral-geometry}

The distributed certificate above proves that a worker presentation is one finite quotient
system. A separate step is needed before observer rows from different shards may be read as a
common neutral geometry.

\begin{definition}[Common quotient chart atlas]\label{def:neutral-quotient-atlas-consensus}
For a shard \(s\), let \(\Sigma_s\) be raw records, let \(\Gamma_s\) be the groupoid of declared
presentation-only moves such as gauge representative changes and local port relabelings, and set
\[
\overline Q_s:=\Sigma_s/\Gamma_s,\qquad X_s:=n_s(\overline Q_s),
\]
where \(n_s\) is the schedule-independent normal-form map. An interface atlas is a family of
domains \(U_{st}\subseteq X_s\), \(U_{ts}\subseteq X_t\), and bijections
\[
\tau_{ts}:U_{st}\to U_{ts}
\]
satisfying identity, inverse, and cocycle laws, with zero closed-path holonomy on graph-shaped
systems. A channel registry assigns each channel \(c\) a metric space \((Z_c,d_c)\), weight
\(w_c>0\), and local features \(F_{s,c}:D_{s,c}\subseteq X_s\to Z_c\). The channel descends
through the atlas when both domain membership and values are transported:
\[
x\in D_{s,c}\Longleftrightarrow \tau_{ts}x\in D_{t,c},
\qquad
F_{t,c}(\tau_{ts}x)=F_{s,c}(x).
\]
\end{definition}

\begin{theorem}[Quotient-visible neutral readout]\label{thm:neutral-quotient-readout-consensus}
Under Definition~\ref{def:neutral-quotient-atlas-consensus}, the relation generated by interface
transport on \(\bigsqcup_sX_s\) is an equivalence relation and defines
\[
Q_{\mathrm{vis}}:=\left(\bigsqcup_sX_s\right)/\!\sim_\tau .
\]
The canonical maps \(X_s\to Q_{\mathrm{vis}}\) agree with transport, and zero closed-path
holonomy makes them injective. Every compatible family of local channel maps descends uniquely
to partial maps \(F_c:Q_{\mathrm{vis}}\dashrightarrow Z_c\). Therefore, on the complete channel
domain
\[
Q_\star=\{x:F_c(x)\ \text{exists for every }c\in\mathcal C_\star\},
\]
the product formula
\[
d_{\mathrm{neu}}(x,y)
=
\left(\sum_{c\in\mathcal C_\star}w_c\,d_c(F_c(x),F_c(y))^p\right)^{1/p}
\]
is a well-defined pseudometric. It is a metric only after quotienting by zero-distance feature
collisions, or after proving that the declared channel family separates points.
\end{theorem}

\begin{proof}
Identity, inverse, and cocycle laws give reflexivity, symmetry, and transitivity. The transport
compatibility of local features is exactly the universal-property condition for descent through
the quotient. The weighted product distance is then independent of the representative. Its
triangle inequality is Minkowski's inequality applied to the channel metrics. The only possible
failure of identity of indiscernibles is equality of all declared feature values, which is
removed by quotienting feature-collision classes or by a joint-separation proof.
\end{proof}

\paragraph{Certificate boundary.}
Pairwise available-channel comparison is not a metric policy: different pairs can share different
channels and violate the triangle inequality. A neutral-geometry run must therefore use complete
cases, a fixed missing symbol whose mask is quotient-visible, or train-only imputation labelled
as an imputed-representation metric. Presentation invariance requires a bijection of
\(Q_{\mathrm{vis}}\) and channel isometries, so gauge changes, port relabelings, observer order,
repair schedule, and shard partition changes are checked on distance matrices, not on displayed
coordinates. Refinement requires a cofinally vanishing tail modulus for the finite-stage
distances. Euclidean claims require the double-centered Gram test
\[
B=-\frac12H(D^{\circ2})H\succeq0
\]
and noisy runs report negative spectral mass, rank/effective rank, held-out stress, and
positive/negative controls. Statistical certificates split independent generative batches before
preprocessing; chart alignment, scaling, imputation, weights, graph construction, dimension
selection, and thresholds are train/validation objects. Any shared shard batch, seed, boundary
condition, trajectory family, duplicate, descendant, or repeated test-set inspection blocks the
held-out theorem. This section certifies a common quotient metric or pseudometric only; physical
Riemannian or Lorentzian spacetime identification belongs to the separate modular/geometric
branch.

\section{Noisy Fair-Block Approximate Consensus}
\label{sec:noisy-fair-block-consensus}

The exact consensus theorem supplies the quotient normal-form target. A noisy implementation
needs one more quantitative certificate: complete fair blocks must contract expected distance to
that target. This section records the conditional bridge from local noisy repair to long-run
observer-facing approximate consensus.

\begin{theorem}[Global noisy approximate consensus under fair-block contraction]
\label{thm:noisy-fair-block-consensus}
Let \((Q,d_Q)\) be the observer-facing quotient state space of a fixed exported OPH patch
federation, and let \(\mathcal N\subset Q\) be the exact quotient normal-form set supplied by the
finite repair theorem on that quotient. Define
\[
D(q):=d_Q(q,\mathcal N)=\inf_{n\in\mathcal N}d_Q(q,n).
\]
Let \(q_{t+1}=\widetilde T_{i_t,t}(q_t)\) be a noisy asynchronous repair process adapted to a
filtration \((\mathcal F_t)_t\). Assume:
\begin{enumerate}[(G1)]
\item \textbf{Exact quotient target.} The ideal repair relation on \(Q\) has the exact OPH
normal-form package: finite exact descent, atomic conflict-component confluence, and repair
completeness, so every fixed initial quotient state has a schedule-independent exact normal form
in \(\mathcal N\).
\item \textbf{Fair asynchronous blocks.} There are stopping times
\[
0=\tau_0<\tau_1<\tau_2<\cdots
\]
such that each interval \([\tau_m,\tau_{m+1})\) contains enough local repairs to expose and act on
every active mismatch class required by the exact transaction/confluence and repair-completeness
certificate, with uniformly bounded length \(\tau_{m+1}-\tau_m\le L\). Write the noisy block map
as
\[
\widetilde B_m
:=
\widetilde T_{i_{\tau_{m+1}-1},\,\tau_{m+1}-1}
\circ\cdots\circ
\widetilde T_{i_{\tau_m},\,\tau_m}.
\]
\item \textbf{Uniform block contraction toward normal form.} There are constants
\(0<\lambda<1\) and \(\varepsilon\ge0\) such that, for every block \(m\) and every reachable
quotient state \(q\) at time \(\tau_m\),
\[
\mathbb E\!\left[
D\!\left(\widetilde B_m(q)\right)
\mid \mathcal F_{\tau_m}
\right]
\le
\lambda D(q)+\varepsilon.
\]
\item \textbf{Controlled within-block excursions.} There are constants \(A\ge1\) and
\(\beta\ge0\) such that, for every \(t\in[\tau_m,\tau_{m+1})\),
\[
\mathbb E\!\left[
D(q_t)\mid \mathcal F_{\tau_m}
\right]
\le
A\,D(q_{\tau_m})+\beta.
\]
\end{enumerate}
Then, at fair-block times,
\[
\mathbb E[D(q_{\tau_m})]
\le
\lambda^mD(q_0)
+
\frac{1-\lambda^m}{1-\lambda}\,\varepsilon,
\]
and hence
\[
\limsup_{m\to\infty}\mathbb E[D(q_{\tau_m})]
\le
\frac{\varepsilon}{1-\lambda}.
\]
At all intermediate asynchronous times,
\[
\limsup_{t\to\infty}\mathbb E[D(q_t)]
\le
A\,\frac{\varepsilon}{1-\lambda}+\beta.
\]
Thus the noisy asynchronous OPH repair process converges in expectation to a controlled tube
around the exact quotient normal-form set. If \(\varepsilon=\beta=0\), then \(D(q_t)\to0\) in
\(L^1\), hence also in probability, along the full asynchronous run.
\end{theorem}

\begin{proof}
Let \(D_m:=D(q_{\tau_m})\). Assumption (G3), applied to the realized state \(q_{\tau_m}\), gives
\[
\mathbb E[D_{m+1}\mid \mathcal F_{\tau_m}]
\le
\lambda D_m+\varepsilon.
\]
Taking expectations,
\[
\mathbb E[D_{m+1}]
\le
\lambda\,\mathbb E[D_m]+\varepsilon.
\]
Iterating the scalar recursion gives
\[
\mathbb E[D_m]
\le
\lambda^mD(q_0)
+
\varepsilon\sum_{r=0}^{m-1}\lambda^r
=
\lambda^mD(q_0)
+
\frac{1-\lambda^m}{1-\lambda}\varepsilon.
\]
Taking \(m\to\infty\) yields the fair-block limsup bound. For
\(t\in[\tau_m,\tau_{m+1})\), Assumption (G4) gives
\[
\mathbb E[D(q_t)]
\le
A\,\mathbb E[D(q_{\tau_m})]+\beta.
\]
Substitution of the fair-block estimate and then taking the limsup over all intermediate times
gives
\[
\limsup_{t\to\infty}\mathbb E[D(q_t)]
\le
A\,\frac{\varepsilon}{1-\lambda}+\beta.
\]
If \(\varepsilon=\beta=0\), then \(\mathbb E[D(q_{\tau_m})]\le\lambda^mD(q_0)\to0\). Since
\(D\ge0\), convergence in expectation to zero implies convergence in probability at block times.
Assumption (G4) then gives \(\mathbb E[D(q_t)]\le A\,\mathbb E[D(q_{\tau_m})]\) inside each
block, so the same \(L^1\) and probability convergence holds along the full asynchronous run.
\end{proof}

\begin{corollary}[Approximate observer-facing schedule independence]
\label{cor:approx-schedule-independence}
Let \(M:Q\to Y\) be an observer-facing readout into a metric space \((Y,d_Y)\), and suppose \(M\)
is \(L_M\)-Lipschitz. Fix an exact sector \(\zeta\) whose exact normal-form set is the singleton
\(\mathcal N_\zeta=\{n_\zeta\}\), and apply Theorem~\ref{thm:noisy-fair-block-consensus} on that
sector, so that \(D(q)=d_Q(q,n_\zeta)\). Then
\[
\limsup_{t\to\infty}
\mathbb E\!\left[
d_Y(M(q_t),M(n_\zeta))
\right]
\le
L_M\left(A\,\frac{\varepsilon}{1-\lambda}+\beta\right).
\]
For two noisy asynchronous schedules \(q_t\) and \(q'_t\) started in the same exact singleton
sector and satisfying the same certificate,
\[
\limsup_{t\to\infty}
\mathbb E\!\left[
d_Y(M(q_t),M(q'_t))
\right]
\le
2L_M\left(A\,\frac{\varepsilon}{1-\lambda}+\beta\right).
\]
\end{corollary}

\begin{proof}
The Lipschitz condition gives
\[
d_Y(M(q_t),M(n_\zeta))
\le
L_M d_Q(q_t,n_\zeta)
=
L_M D(q_t).
\]
The first bound follows from Theorem~\ref{thm:noisy-fair-block-consensus}. The two-schedule bound
follows from the triangle inequality through \(M(n_\zeta)\) and applying the first estimate to
both schedules.
\end{proof}

\begin{corollary}[High-probability noisy consensus tube]
\label{cor:high-prob-noisy-tube}
Assume the block-distance process also admits the pathwise decomposition
\[
D_{m+1}
\le
\lambda D_m+\varepsilon+\xi_{m+1},
\qquad
\mathbb E[\xi_{m+1}\mid\mathcal F_{\tau_m}]=0,
\qquad
|\xi_{m+1}|\le b
\]
almost surely. Then, for every \(a>0\),
\[
\Pr\!\left[
D_m>
\lambda^mD_0
+
\frac{1-\lambda^m}{1-\lambda}\varepsilon
+
a
\right]
\le
\exp\!\left(
-\frac{a^2(1-\lambda^2)}{2b^2}
\right).
\]
\end{corollary}

\begin{proof}
Unrolling the recursion gives
\[
D_m
\le
\lambda^mD_0
+
\frac{1-\lambda^m}{1-\lambda}\varepsilon
+
\sum_{r=1}^{m}\lambda^{m-r}\xi_r.
\]
The final term is a weighted martingale sum with increments bounded by
\(|\lambda^{m-r}\xi_r|\le\lambda^{m-r}b\). Azuma--Hoeffding gives
\[
\Pr\!\left[
\sum_{r=1}^{m}\lambda^{m-r}\xi_r>a
\right]
\le
\exp\!\left(
-\frac{a^2}{2\sum_{r=1}^{m}\lambda^{2(m-r)}b^2}
\right).
\]
Since \(\sum_{r=1}^{m}\lambda^{2(m-r)}\le(1-\lambda^2)^{-1}\), substitution yields the claimed
bound.
\end{proof}

\begin{proposition}[Finite fair-block contraction certificate]
\label{prop:finite-fair-block-certificate}
Let \(Q\) be finite, let \(\mathfrak B_{\mathrm{fair}}\) be a finite list of noisy fair-block
types, and let \(K_B(q,q')\) be the Markov kernel induced by block type \(B\). If there are
constants \(0<\lambda<1\) and \(\varepsilon\ge0\) such that
\[
\sum_{q'\in Q}K_B(q,q')D(q')
\le
\lambda D(q)+\varepsilon
\qquad
\forall q\in Q,\quad \forall B\in\mathfrak B_{\mathrm{fair}},
\]
then Assumption (G3) of Theorem~\ref{thm:noisy-fair-block-consensus} holds for any run whose
fair blocks are drawn from \(\mathfrak B_{\mathrm{fair}}\).
\end{proposition}

\begin{proof}
Conditioning on the current quotient state \(q\) and the realized fair-block type \(B\), the
conditional expectation of \(D\) after the block is exactly
\(\sum_{q'}K_B(q,q')D(q')\). The displayed inequality is therefore Assumption (G3), uniformly
over all reachable states and fair-block types.
\end{proof}

\paragraph{Finite audit route and constants.}
For a finite exported packet net, Proposition~\ref{prop:finite-fair-block-certificate} gives the
practical certificate path:
\[
\text{finite quotient }Q
\Rightarrow
\text{exact normal forms }\mathcal N
\Rightarrow
\text{distance table }D(q)
\Rightarrow
\text{noisy fair-block kernels }K_B
\Rightarrow
(\lambda,\varepsilon)\text{ certificate}.
\]
Here \(\mathcal N\) is the exact quotient normal-form set, \(D(q)\) is observer-facing residual
distance to that set, \(\lambda\) is the net contraction produced by one completed fair block,
\(\varepsilon\) is the per-block irreducible local recovery / record / readout / calibration /
environmental noise, \(A\) bounds transient within-block expansion, \(\beta\) is the
within-block noise floor, and \(L\) is the fairness horizon before all required active mismatch
classes are serviced. In quantum/collar implementations, \(\varepsilon\) may absorb Petz-domain
truncation, Fawzi--Renner recovery error, approximate central-record error, detector/readout
noise, and coarse-graining defect.

\paragraph{Claim boundary.}
Theorem~\ref{thm:noisy-fair-block-consensus} is a conditional global noisy-consensus theorem. It
does not follow from the exact finite repair theorem alone. The exact OPH theorem supplies the
quotient normal-form target \(\mathcal N\); the noisy theorem requires a separate fair-block
contraction certificate \((\lambda,\varepsilon,A,\beta,L)\) for the chosen implementation. Without
that certificate, OPH retains exact fixed-cutoff convergence and the collar-local splice and
record-stability estimates above. With that certificate, arbitrarily long asynchronous noisy
repair sequences converge to a controlled observer-facing tube around the exact quotient
normal-form set.


\section{Law-Space Selection and Observer Emergence}

This section studies a simple meta-selection model on law space. The aim is to formalize one
criterion for favoring schedule-robust, observer-supporting, and simple laws; the replicator
dynamics below is part of the model, not a claim about literal cosmological dynamics.

We begin by defining what it means for a law to support observers. An observer is treated
operationally as a \textbf{persistent predictive module}: a subgraph that maintains a stable
record algebra and uses its output law to predict its boundary's future behavior.

\begin{definition}[Schedule robustness]\label{def:robustness}
Fix distributions $\mu$ over initial conditions and $\nu$ over asynchronous schedules, and a gauge-invariant observable $M$. For a law $\lambda$, define
\[
\mathcal{R}_M(\lambda)
:=
\Pr_{s\sim\mu,\;\sigma,\tau\sim\nu}
\!\left[
M\bigl(\operatorname{nf}^{\sigma}_\lambda(s)\bigr)
=
M\bigl(\operatorname{nf}^{\tau}_\lambda(s)\bigr)
\right].
\]
If Theorem~\ref{thm:confluence} holds for $\lambda$, then $\mathcal{R}_M(\lambda)=1$.
\end{definition}

\begin{definition}[Observer yield]\label{def:yield}
Let $X_t^\lambda$ denote the stationary process obtained by repeated local perturbation plus
reconciliation under law $\lambda$. For each subgraph $U\subseteq V$, let
\(\mathcal Z_U^{\mathrm{rec}}(t)\) be the declared exact record algebra on its observer-accessible
surface, or the reference exact algebra when only an approximate record presentation is available,
and let $Y_U(t)$ be the corresponding finite outcome variable induced by that record algebra. Then $U$ is
\((\eta,\varepsilon,h)\)-observer-like if it is record-stable:
\[
d_{\mathrm{TV}}\!\bigl(\operatorname{Law}(Y_U(t+h)),\operatorname{Law}(Y_U(t))\bigr)\le \eta,
\]
and predictive:
\[
I\bigl(Y_U(t);\; X_{\partial U,\,t+1:t+h}^\lambda\bigr)\ge\varepsilon.
\]
Define
\[
\mathcal{O}_{\eta,\varepsilon,h}(\lambda)
:=
\mathbb{E}\!\left[
\#\left\{
U\subseteq V:
U \text{ is } (\eta,\varepsilon,h)\text{-observer-like}
\right\}
\right].
\]
\end{definition}

\begin{definition}[Law fitness]\label{def:fitness}
Let $K(\lambda)$ be a description-length penalty. Define
\[
f(\lambda) = \alpha\,\mathcal{R}_M(\lambda) + \beta\,\mathcal{O}_{\eta,\varepsilon,h}(\lambda) - \gamma\,K(\lambda),
\]
with $\alpha,\beta,\gamma>0$.
\end{definition}

\begin{theorem}[Replicator monotonicity on law space]\label{thm:replicator}
Let $\Lambda=\{\lambda_1,\dots,\lambda_m\}$ be candidate laws with population weights $x_i(t)$ under replicator dynamics:
\[
\dot{x}_i = x_i(f_i - \bar{f}),
\qquad
f_i := f(\lambda_i),
\qquad
\bar{f} := \sum_{j=1}^m x_j f_j.
\]
Then
\[
\frac{d}{dt}\bar{f} = \operatorname{Var}_x(f) \ge 0.
\]
Mean fitness is nondecreasing, and strictly increasing unless all extant laws have the same fitness.
\end{theorem}

\begin{proof}
Direct computation:
\[
\frac{d}{dt}\bar{f}
=
\sum_i \dot{x}_i f_i
=
\sum_i x_i(f_i - \bar{f})f_i
=
\sum_i x_i f_i^2 - \bar{f}^2
=
\operatorname{Var}_x(f) \ge 0.
\]
Equality iff all $f_i$ on the support of $x$ are equal.
\end{proof}

This theorem records the monotonicity property of the meta-selection model.


\section{Connection to Observer-Patch Holography}

The formalism above is the computational skeleton of Observer-Patch Holography (OPH). The
observer patches carry von Neumann algebras on support-visible holographic cuts. In symmetric
regulator charts those cuts may be represented by patches on a screen \(S^2\). The fixed-cutoff
microphysics carrier is a federated patch system with echosahedral local interfaces. The \(S^2\)
chart supplies cap and collar geometry, and in the companion relativity branch its conformal group
supplies the Lorentz bridge. The carrier supplies finite ports, records, and repair interfaces. The
overlap projections are restrictions to shared subalgebras, and the consistency condition is
algebraic state agreement on overlaps.

The bridge to physics works as follows in the broader companion corpus:

\begin{itemize}
\item The patch net becomes a net of support-visible subregion algebras; \(S^2\) is the standard
observer-facing regulator chart, not a required literal substrate.
\item Overlap Consistency, one of the canonical framework axioms, is the algebraic version of Definition~\ref{def:patchnet}.
\item The recoverability clause of the canonical Recoverable Generalized Entropy axiom provides controlled collar recovery data. The exact collar factorization \(\rho_{ABD} = \bigoplus_\alpha p_\alpha\,\rho^{(\alpha)}_{Ab_L^\alpha}\otimes\rho^{(\alpha)}_{b_R^\alpha D}\) is used only at exact Markovity or along the fixed-collar replacement limit, while the declared Petz/Fawzi--Renner recovery channels supply recovered comparison states from which the local repair moves are built (see Appendix~\ref{app:splice} and Definition~\ref{def:oph-repair-map}).
\item Gauge symmetry as implementation hiding (Theorem~\ref{thm:gauge}) becomes the fixed-cutoff edge-sector package from which the compact-gauge branch, using the theorem-produced coherent refinement transport and compatible finite fibers, reconstructs a compact gauge group from zero-obstruction sectors. Zero obstruction is the transportability/classification condition; it is not the Standard Model selector. The realized witness theorem supplies nonempty MAR-admissible branch data, and MAR plus the explicit one-Higgs matter package then selects the Standard Model quotient $\mathrm{SU}(3)\times\mathrm{SU}(2)\times\mathrm{U}(1)/\mathbb{Z}_6$ together with the exact hypercharge lattice and the realized counts \(N_c=3\), \(N_g=3\).
\item On the declared support-visible compact-gauge branch, the companion compact paper obtains
the four-dimensional Euclidean Yang--Mills form from compact-gauge holonomy data and the local
MaxEnt/Gibbs continuum limit, with the branch assumptions stated there. It then applies the same
quotient-first repair logic to Euclidean transfer: active exact-Markov repairs become conditional
expectations, their finite-stage generator has a uniform repair gap, and coherent continuum
extraction identifies the Yang--Mills gap with the repair gap. The Clay-facing admissibility claim
is exactly that this branch-local support-visible extraction and its
Osterwalder--Schrader reconstruction supply the required four-dimensional axiomatic construction.
\item The coarse-graining compatibility theorem (Theorem~\ref{thm:coarse-grain-consensus}) is the link between the finite reconciliation protocol and the refinement/RG language used by the OPH branches: the macroscopic law space is stable when the selected coarse-graining channel shadows the normal-form and obstruction maps with controlled defects.
\item Stable defects (Corollary~\ref{cor:defects}) become the topologically protected excitations identified with particles on the declared branch.
\item The record-algebra theorem (Theorem~\ref{thm:crdt}) provides the formal basis for the fixed-cutoff observation layer, where records are carried by central or quantitatively stable approximately commuting projectors in overlap centers.
\end{itemize}

The companion manuscripts develop a derived gravity branch from entanglement equilibrium and modular geometry, the SM gauge-group and count closure from edge-sector reconstruction plus the realized MAR branch, a support-visible compact-gauge Yang--Mills form and repair-gap theorem under the declared branch assumptions, and a controlled large-\(N_{\mathrm{edge}}\) worldsheet effective-description branch above the heat-kernel edge-sector identity. The implemented cosmological capacity branch, with its zero-input self-closure target, and later phenomenological continuations are kept explicitly separate from that recovered-core claim set.

This paper provides the finite patch-net foundation for the broader companion corpus.


\section{Discussion and Scope Boundaries}

This paper proves the fixed-point consensus spine of OPH: schedule-independent normal forms from
fixed initial quotient states, boundary-conditioned uniqueness when a preserved boundary/sector
fiber has a unique consistent extension, nontrivial functional selected-fiber branch elimination,
holonomy obstructions, gauge-quotient invariance,
separated cofinal refinement-limit consensus classes, controlled coarse-graining compatibility,
the fixed-cutoff operator-record theorem, distributed one-universe realization for admissible
worker presentations of a single global carrier, and conditional noisy fair-block convergence once a
separate contraction certificate is supplied.
It also proves the full repair-completeness, Petz-domain, and quotient-compatibility package on
the rooted-tree packet-net domain of Theorem~\ref{thm:tree-packet-domain}, and gives a clean
law-selection meta-model. The relativity chain and the realized Standard Model
structural chain are the recovered core. The capacity relation is a separate implemented branch
with a zero-input self-closure target. Downstream phenomenology requires additional assumptions
beyond the consensus results proved here.

\paragraph{Complexity boundary.}
On a fixed finite patch net, the accepted reconciliation dynamics is a finite-state asynchronous
rewrite system on \(\Sigma\). Under Theorem~\ref{thm:confluence}, every accepted repair run has at
most the number of distinct reachable \(\mu\)-values below its start value minus one, because each
accepted transaction strictly lowers the declared exact measure. On the scalar \(\Phi\)-branch
this specializes to the older \(|\Phi(\Sigma)|-1\le |\Sigma|-1\) bound. Exact normal-form
computation is therefore decidable by direct iteration of accepted aggregate transactions. This
step bound is a termination statement; the schedule-independent answer depends on the atomic
conflict-component local diamond and repair-completeness clauses of Theorem~\ref{thm:confluence}.
What this paper does
\emph{not} prove is a uniform polynomial-time bound, a sharper complexity-class placement for
families of growing patch nets, or any hardness lower bound.

\paragraph{Approximate-stability boundary.}
The theorem-grade consensus statement is exact on the declared fixed-cutoff branch. Approximate
control begins collar-locally: Theorem~\ref{thm:splice} gives the exact-Markov modulus
\(\delta^{\mathrm M}_{A:B:D}(\varepsilon)\to0\) on one fixed finite-dimensional collar model and
the one-shot recovery comparison bound \(2\sqrt{1-e^{-\varepsilon}}\le 2\sqrt{\varepsilon}\),
while Theorem~\ref{thm:crdt} gives the \((\varepsilon,\delta_{\mathrm{rec}})\) repeated-read
stability bound for approximate record projectors. The long-run noisy statement is conditional:
Theorem~\ref{thm:noisy-fair-block-consensus} upgrades those local noisy controls to a global
expected tube around the exact quotient normal-form set only after a fair-block contraction
certificate \((\lambda,\varepsilon,A,\beta,L)\) is supplied for the chosen implementation. It does
not follow from finite descent or fairness alone, and it gives a unique approximate readout only
inside singleton boundary/sector fibers.

\paragraph{Expressive-power boundary.}
The law-selection model of Theorem~\ref{thm:replicator} is a finite-candidate monotonicity result.
For each fixed patch net the theorem package proves a finite-state exact reconciliation
mechanism. A universality claim would require an explicit uniform family of patch nets and repair laws that
simulates arbitrary circuits or machines with stated encoding overhead and robustness under
asynchronous schedules. No such theorem is supplied here.

With those boundaries explicit, the consensus-paper boundaries and companion interfaces are:

\begin{enumerate}
\item \textbf{Sharper RG-shadowing estimates.} Theorem~\ref{thm:coarse-grain-consensus} closes the abstract reconciliation/coarse-graining square once its normal-form and obstruction defects are supplied. The quantitative task is to derive model-specific or uniform bounds for \(\varepsilon^n_{sr}\) and \(\varepsilon^h_{sr}\) from concrete OPH coarse-graining channels and recovery decoders.

\item \textbf{Defect classification and refinement-limit transportability.} The fixed-cutoff hierarchy extends from abelian frustrations to crossed-module classes \(q\in \check H^2(N,H\to G)\). The continuation task is to connect those higher-gauge defect sectors to the refinement-stable transportable sector category used in the broader compact-gauge reconstruction lane.

\item \textbf{Observable-level confluence beyond the declared fixed-cutoff physical algebra.} Theorem~\ref{thm:obsconfluence} closes the fixed-cutoff quantum-lift statement when microscopic representatives differ by gauge relabelings globally or by sector/higher-gauge relabelings on the same declared quotient-local glued state. The continuation question is whether an analogous observable theorem survives on broader refinement-stable branches where the union-collar compatibility is only approximate or where the physical observable algebra itself changes under refinement.

\item \textbf{Distributed implementation certificates.} Theorem~\ref{thm:distributed-one-universe} proves invariance for a worker presentation only after the implementation supplies one global carrier, authoritative ownership, cut interfaces, projection-preserving events, rollback roots, and final monolithic normal-form/readout certificates. Shard-local seeds, missing cut artifacts, stale manifests, synthetic seam trajectories, or worker-id histories do not certify a one-universe realization.

\item \textbf{Global approximate-consensus stability.} Theorem~\ref{thm:splice} and Theorem~\ref{thm:crdt} supply the collar-local perturbative controls carried by this paper. The rooted-tree packet domain proves the exact finite repair package on one nontrivial exported domain. Theorem~\ref{thm:noisy-fair-block-consensus} closes the long-run noisy branch only after a fair-block contraction certificate is supplied. What remains absent is an automatic theorem showing that arbitrary approximate recovery moves on arbitrary packet-closed exported overlap nets are repair-complete, preserve transactional validation and atomic conflict-component confluence, satisfy the support/CPTP clause on every Petz branch used there, and meet that fair-block contraction certificate.

\item \textbf{Expressive power / universality.} A universality claim would require an explicit uniform family of patch nets and repair laws that simulates arbitrary circuits or machines with stated encoding overhead and robustness under asynchronous schedules. This paper supplies no such construction, so universality is outside the theorem-grade output.

\item \textbf{Interface to the companion gravity/gauge stack.} The companion compact paper derives the gravity chain through the support-visible BW scaling theorem on the prime geometric cap subnet. This consensus paper supplies the finite-state and refinement-consensus spine used by that bridge, while dark-sector proposals, heuristic baryogenesis continuations, spectroscopy, and string/worldsheet topics remain separate continuations that require their own declared inputs beyond what this paper itself proves.
\end{enumerate}

The last item is an interface statement rather than an extra consensus premise. The fixed-point
theorems proved here stand on their stated finite and refinement-consensus hypotheses, while
companion gravity, gauge, and continuation sectors add their own theorem surfaces and declared
inputs.

\subsection{Assumption-Dependent BFT and QECC Extensions}
\label{subsec:assumption-dependent-bft-qecc}

The consensus formalism of \textrm{OPH} has natural analogies to classical and quantum
distributed Byzantine agreement. Observer patches correspond to protocol nodes, overlap repair
corresponds to a quorum vote, and the repair fixed-point corresponds to a consensus state. Under
explicit structural assumptions (quorum size $\geq 2f+1$, partial synchrony, one-vote-per-view,
certificate semantics, and \textrm{DLS}-style view-change), a \textrm{QBFT}-style interpretation
of \textrm{OPH} repair satisfies safety and liveness (Appendix~\ref{app:bft-qecc},
Theorem~\ref{thm:qbft-safety}). On the fixed-cutoff collar branch used here, the repair map is
written in exact-splice / Petz form; the assumption-dependent item is the \textrm{CPTP}
property on all inputs, which requires either full-rank $\mathcal{N}(\sigma)$ or an explicit
domain restriction, and trace-preserving completion is not automatic when
$\mathcal{N}(\sigma)$ has a non-trivial kernel
(Proposition~\ref{prop:petz-cptp}). A quantum error-correcting interpretation is possible only
after a genuine code subspace, logical dictionary, error family, and recovery map are supplied.
The graph-min-cut equality for distance is not a property of a bare overlap graph: the same graph
can realize distance \(1\) or distance \(|V|\) under different interface maps
(Theorem~\ref{thm:no-free-mincut}). Distance/min-cut statements require the topological-code
certificate of Definition~\ref{def:tcert} and Theorem~\ref{thm:tcert-distance}; resilience
requires the Knill--Laflamme certificate of Theorem~\ref{thm:qecc}. All of these extensions are
assumption-dependent or conjectural and are not part of the core theorem package of Paper~4.

\begin{center}
\small
\begin{tabular}{@{}p{0.33\linewidth}p{0.56\linewidth}@{}}
\textbf{Desired statement} & \textbf{Required certificate} \\
\hline
Overlap network is a code & finite constraint-code data of Proposition~\ref{prop:constraint-code} \\
Repair converges & finite exact descent; confluence additionally needs atomic conflict-component local diamond and repair completeness \\
Distance equals min-cut & topological-code certificate with homological logicals and matching systole/min-cut geometry \\
Corrects \(t\) corrupted carriers & code projector, error family, Knill--Laflamme condition, and certified \(t<d/2\) distance bound \\
Exponential convergence & declared transfer/channel operator with stationary projection and spectral gap \\
Long-run noisy approximate consensus & fair-block contraction certificate \((\lambda,\varepsilon,A,\beta,L)\) for distance to the exact quotient normal-form set \\
Wall-clock BFT liveness & partial synchrony, quorum certificates, authentication, and \(n\ge3f+1\) \\
Hardware search work reduction & exact-verifier candidate-enrichment factor measured under controls \\
\end{tabular}
\end{center}


\appendix
\section{Quantum/Algebraic Lift: Markov-Collar Splice Theorem}\label{app:splice}

This appendix records the algebraic splice statement used to relate the finite patch-net model to
the OPH collar formalism.

An observer is written as
\[
O=(P,\mathcal{A}(P),\rho,R),
\]
where $P$ is the screen patch, $\mathcal{A}(P)$ the local von Neumann algebra, $\rho$ the local state, and $R$ the record algebra.

\begin{theorem}[Markov-collar splice theorem, exact and controlled]\label{thm:splice}
Suppose a collar tripartition $A$-$B$-$D$ has exact Markov decomposition
\[
\rho_{ABD}
=
\bigoplus_{\alpha}
p_\alpha\,
\rho^{(\alpha)}_{A b_L^\alpha}
\otimes
\rho^{(\alpha)}_{b_R^\alpha D}.
\]
Let $\sigma_{b_R^\alpha D'}^{(\alpha)}$ be any family of normalized environment states compatible with the same right-boundary sectors. Define
\[
\rho'_{AB D'}
=
\bigoplus_{\alpha}
p_\alpha\,
\rho^{(\alpha)}_{A b_L^\alpha}
\otimes
\sigma_{b_R^\alpha D'}^{(\alpha)}.
\]
Then for every observable $X$ supported on $A\cup b_L$,
\[
\operatorname{Tr}(X\rho'_{AB D'})
=
\operatorname{Tr}(X\rho_{ABD}).
\]
Fix one finite-dimensional collar model and let
\[
\mathfrak M_{A:B:D}
:=
\left\{
\tau_{ABD}: I(A:D\mid B)_\tau=0
\right\},
\]
with exact-Markov distance modulus
\[
\delta^{\mathrm M}_{A:B:D}(\varepsilon)
:=
\sup\left\{
\inf_{\tau\in\mathfrak M_{A:B:D}}\|\omega-\tau\|_1:
I(A:D\mid B)_\omega\le\varepsilon
\right\}.
\]
Then
\[
\delta^{\mathrm M}_{A:B:D}(\varepsilon)\to0
\qquad
(\varepsilon\downarrow0).
\]
On a fixed faithful collar class with lower floor \(\lambda_\ast>0\), this
qualitative modulus can be sharpened to a collar-local rate
\[
\delta^{\mathrm M,\lambda_\ast}_{A:B:D}(\varepsilon)
\le
C_{A:B:D,\lambda_\ast}\,\varepsilon^{\theta_{A:B:D,\lambda_\ast}},
\]
by the compact real-analytic Lojasiewicz inequality applied to
\(I(A:D\mid B)\). The constants depend on the fixed collar model and floor;
they are not a dimension-free stability theorem for arbitrary tripartite
systems.
Hence if \(I(A:D\mid B)_\omega\le \varepsilon\) and \(\widetilde\omega_\varepsilon\in\mathfrak M_{A:B:D}\) is chosen so that
\[
\|\omega-\widetilde\omega_\varepsilon\|_1
\le
\delta^{\mathrm M}_{A:B:D}(\varepsilon),
\]
the corresponding exact splice \(\widetilde\omega'_\varepsilon\) satisfies
\[
\left|
\operatorname{Tr}(X\omega)-\operatorname{Tr}(X\widetilde\omega'_\varepsilon)
\right|
\le
\|X\|_\infty\,
\delta^{\mathrm M}_{A:B:D}(\varepsilon)
\]
for every observable \(X\) supported on \(A\cup b_L\).

Independently, if \(I(A:D\mid B)_\omega\le \varepsilon\), then there exists a recovery map \(\mathcal R_{B\to BD}\) such that
\[
\left\|
\omega_{ABD}
-
(\mathrm{id}_A\otimes \mathcal R_{B\to BD})(\omega_{AB})
\right\|_1
\le
2\sqrt{1-e^{-\varepsilon}}
\le
2\sqrt{\varepsilon}.
\]
\end{theorem}

\begin{proof}
The exact splice statement is the usual blockwise factorization argument:
\[
\operatorname{Tr}(X\rho'_{AB D'})
=
\sum_\alpha
p_\alpha\,
\operatorname{Tr}\!\left(
X\, \rho^{(\alpha)}_{A b_L^\alpha}
\right)
\operatorname{Tr}\!\left(\sigma_{b_R^\alpha D'}^{(\alpha)}\right).
\]
Each right factor is normalized, so the value agrees with the same computation for \(\rho_{ABD}\).

For the controlled statement, compactness of the fixed finite-dimensional state space and continuity of conditional mutual information imply \(\delta^{\mathrm M}_{A:B:D}(\varepsilon)\to0\): otherwise one could find a sequence with \(I(A:D\mid B)\to0\) staying a fixed trace distance away from every exact Markov state, contradicting convergence of a subsequence to an exact Markov limit point. Once \(\widetilde\omega_\varepsilon\) is chosen, the exact splice identity for \(\widetilde\omega_\varepsilon\) gives
\[
\left|
\operatorname{Tr}(X\omega)-\operatorname{Tr}(X\widetilde\omega'_\varepsilon)
\right|
=
\left|
\operatorname{Tr}\!\left[X(\omega-\widetilde\omega_\varepsilon)\right]
\right|
\le
\|X\|_\infty\,
\delta^{\mathrm M}_{A:B:D}(\varepsilon).
\]
The final inequality is the standard Fawzi--Renner recovery bound~\cite{fawziRenner15}.
\end{proof}

This appendix therefore uses exact splice identities in only two regimes: literal exact Markovity, or a controlled collar family on one fixed finite-dimensional model for which \(\delta^{\mathrm M}_{A:B:D}(\varepsilon)\to0\). Fawzi--Renner recovery supplies the constructive recovered comparison state; the fixed-collar modulus supplies the exact-Markov comparison. Small one-shot conditional mutual information is not silently upgraded to an exact normal form.

\section{Fixed-Cutoff Realization, Quotient Repair, and Edge Centers}
\label{app:consensus-fixed-cutoff-port}

This appendix carries the fixed-cutoff realization and edge-center items for the consensus paper.
They sharpen the
quotient-first repair interpretation used throughout the consensus paper and make the collar
boundary data explicit at the same finite patch-net level. On the declared fixed-cutoff collar
branch, the local repair step is read from exact Markov splice or a declared Petz/Fawzi--Renner
recovery move on that same collar data; the recovery move gives a recovered comparison state, while
exact splice requires exact Markovity or a controlled fixed-collar replacement modulus. Representative
repair maps are only lifts of the resulting quotient-local update.

\subsection{Quotient Repair and UV Underdetermination}

At fixed cutoff, each regulator cell \(x\) carries a finite-dimensional factor
\(\mathfrak h_x\), patch algebras are finite type-I algebras, and gauge-as-gluing is realized as
a compact boundary redundancy action on cut data. The physical repair law therefore belongs on the
overlap-invariant quotient rather than on hidden representatives. If \(q:\Sigma\to\Sigma/\Gamma\)
is the quotient by boundary redundancy and \(\overline T_i\) is the physical quotient update, a
representative-level map \(T_i\) is only required to be a lift satisfying
\[
q\circ T_i=\overline T_i\circ q.
\]
Hence
\[
q(T_i(\gamma\cdot s))=q(T_i(s))
\]
for gauge-equivalent inputs. Quotient descent is therefore structural, while strict
representative-level covariance is only implementation bookkeeping. The burden is to prove that
the accepted recovery-derived local moves satisfy the stated repair-completeness, support-local
disjoint-commutation, nested-collar
restriction-compatibility, and Petz-domain control clauses on the declared branch. The
touched-overlap acceptance contract yields finite Lyapunov descent and derived
termination for accepted moves, while the fixed-cutoff gluing package carries the
parenthesization-invariant union-collar state used for the local diamond.

\begin{proposition}[Ancilla-stable UV underdetermination]
Let a fixed-cutoff OPH realization be stabilized by finite ancillary factors \(K_P\) in a fixed
product state, with observable patch algebras embedded as
\(\mathcal A(P)\otimes \mathbf 1_{K_P}\) and repair dynamics acting trivially on the ancillas.
Then observable expectations on the physical subalgebras, overlap data, the local-Gibbs branch,
the collar conditional mutual information \(I(A:D\mid B)\), the Fawzi--Renner remainder, the
collar Markov modulus, and the quotient normal form are unchanged. Thus the fixed-cutoff theorem
package determines the UV branch only modulo such ancillary stabilization together with gauge or
implementation hiding, not a unique microscopic presentation.
\end{proposition}

\begin{proof}
Product ancillas leave physical observables unchanged, cancel additively inside conditional mutual
information, and are inert under the repair maps. Hence every invariant listed above is unchanged.
\end{proof}

\subsection{Derived Boundary Data and Ordinary EC}

\begin{proposition}[Derived boundary gluing datum]
Choose a finite regulator chart for the patches meeting along a connected cut \(\Sigma\). Because
the local overlap algebras are finite-dimensional matrix algebras, any overlap-consistent
recharting is an inner automorphism and is implemented by a unitary on the cut Hilbert space. The
compact closure of the subgroup generated by these recharting unitaries is a compact boundary
redundancy group \(K_\Sigma\). If triple-overlap defects are central, the projective composition
law lifts to a compact central extension \(\widehat K_\Sigma\); on the ordinary branch one simply
sets \(\widehat K_\Sigma = K_\Sigma\). A genuinely noncentral \(2\)-group defect is the only
obstruction to reducing the overlap transition system to an ordinary compact group action.
\end{proposition}

\begin{theorem}[Derived EC decomposition]
Under the fixed-cutoff regulator realization above, and on the ordinary or central-defect branch,
the collar Hilbert space is
\[
\mathcal H_{B_\delta}
=
(\tilde{\mathcal H}_{B_L}\otimes \tilde{\mathcal H}_{B_R})^{\widehat K_\Sigma}
\cong
\bigoplus_{\alpha}
\left(\mathcal H_{b_L^\alpha}\otimes \mathcal H_{b_R^\alpha}\right),
\]
and the center of the collar algebra is generated by the block projectors:
\[
Z(\mathcal A(B_\delta))
=
\bigoplus_\alpha \mathbb C\cdot \mathbf 1_\alpha.
\]
The right half-collar carries the contragredient representation because it sees inverse transport
across the same cut.
\end{theorem}

\begin{remark}
This is the finite-patch-net origin of the collar center used by the later Markov, record, and
observer packages. Exact Markovity is an additional state hypothesis; EC provides the
kinematic block structure.
\end{remark}

\subsection{Higher-Gauge Replacement on the Genuinely Noncentral Branch}

\begin{proposition}[Derived higher-gauge cut datum]
On the genuinely noncentral branch, weak overlap gluing on a connected cut \(\Sigma\) is encoded
by a compact crossed module
\[
\mathbb K_\Sigma=(H_\Sigma\xrightarrow{\partial_\Sigma}G_\Sigma,\triangleright)
\]
with defect class
\[
q_\Sigma\in \check H^2(N_\Sigma,H_\Sigma\to G_\Sigma),
\]
and compact higher-gauge change system
\[
\mathcal T_\Sigma=C^1(N_\Sigma,H_\Sigma)\rtimes C^0(N_\Sigma,G_\Sigma).
\]
\end{proposition}

\begin{theorem}[Higher-gauge EC decomposition and defect transport]
On the genuinely noncentral branch,
\[
\mathcal H_{B_\delta}^{2g}
=
(\tilde{\mathcal H}_{B_L}\otimes \tilde{\mathcal H}_{B_R})^{\mathcal T_\Sigma}
\cong
\bigoplus_\lambda
(\mathcal H_{b_L^\lambda}\otimes \mathcal H_{b_R^\lambda}),
\]
and
\[
Z(\mathcal A_{2g}(B_\delta))
=
\bigoplus_\lambda \mathbb C\cdot \mathbf 1_\lambda.
\]
Moreover the full crossed-module orbit \(q_\Sigma\) is invariant under local rechartings and
classifies fixed-cutoff genuinely noncentral gluing data. The higher associator is removable iff
\(q_\Sigma\) lies in the image of
\[
\check H^1(N_\Sigma,G_\Sigma)longrightarrow
\check H^2(N_\Sigma,H_\Sigma\to G_\Sigma),
\qquad [g]\longmapsto[(g,1)].
\]
That map need not be injective. Strict endpoint-only ordinary transport additionally requires at
least one allowed strict representative with trivial represented loop holonomy.
\end{theorem}

\begin{corollary}[Exact Markov plus split alignment adds the state factorization]
On either the ordinary/central branch of the previous theorem or the genuinely noncentral
higher-gauge branch, if in addition
\[
I_\omega(A_\delta:D_\delta\mid B_\delta)=0
\]
\emph{and} the state satisfies the Markov-split alignment hypothesis (its HJPW decomposition of
\(\mathcal H_{B_\delta}\) can be chosen to be the EC decomposition itself; see the main-text
Section 2.3 and the compact paper's alignment definition), or one passes to the explicitly stated
idealized recoverability limit that supplies both conditions, then
\[
\rho_{A_\delta B_\delta D_\delta}
=
\bigoplus_\alpha p_\alpha
\left(\rho_{A_\delta b_L^\alpha}\otimes \rho_{b_R^\alpha D_\delta}\right).
\]
EC therefore gives the kinematic block decomposition, while exact Markovity plus split alignment
is the extra state input that gives the EC-aligned HJPW normal form. Exact Markovity alone yields
the normal form only over a state-dependent HJPW split, which a Bell-pair example (main text,
Section 2.3) shows can be transposed relative to the EC factors.
\end{corollary}

\begin{proposition}[Parenthesization-invariant union-collar gluing]\label{prop:assocglue}
Fix a finite union collar \(U\) built from two overlapping local repair collars on the declared
fixed-cutoff branch. On the ordinary or central-defect branch, the physical glued state on \(U\)
determined by the exterior marginals and sector data is independent of how the local gluings are
parenthesized, up to the boundary-redundancy action. On the genuinely noncentral branch, the same
statement holds up to the crossed-module change system \(\mathcal T_\Sigma\). Hence the physical
quotient-local glued state on \(U\) is well defined independently of parenthesization.
\end{proposition}

\begin{proof}
On the ordinary or central-defect branch, the preceding edge-center decomposition writes the
collar algebra as a direct sum of sector blocks generated by central projectors. Reparenthesizing
two overlapping local gluings changes only the representative by the central boundary action and
cannot move the state between those block projectors, so the quotient-local glued state is
unchanged. On the genuinely noncentral branch, the failure of strict composition is recorded by
the crossed-module associator, and the higher-gauge theorem above identifies rechartings exactly
with the \(\mathcal T_\Sigma\)-coboundary action. Reparenthesization therefore changes only the
representative inside one \(\mathcal T_\Sigma\)-orbit. In either case the physical quotient state
is parenthesization-independent.
\end{proof}

\begin{remark}
Approximate recoverability gives controlled deviations from this normal form; it is not implied by
EC alone. This is the fixed-cutoff topological package behind the consensus paper's quotient and
record-language surface.
\end{remark}

\begin{corollary}[Physical observables are invariant on one quotient-local glued state]\label{cor:obsglue}
Let \(U\) be a finite union collar on the declared fixed-cutoff branch, and let
\(\omega_U,\omega'_U\) be two microscopic representatives of the same quotient-local glued state
from Proposition~\ref{prop:assocglue}. Then every physical observable \(X\) on the collar
fixed-point / quotient-local algebra has the same expectation in both representatives:
\[
\operatorname{Tr}(X\omega_U)=\operatorname{Tr}(X\omega'_U).
\]
In particular the same holds for the central sector projectors and for any observer-accessible
record observable generated from them on that same declared surface.
\end{corollary}

\begin{proof}
On the ordinary or central-defect branch, Proposition~\ref{prop:assocglue} says the two
representatives differ only by the boundary-redundancy action inside one fixed sector block. The
fixed-point collar algebra and its central block projectors are invariant under that action, so
their expectation values agree. On the genuinely noncentral branch, the same proposition says the
two representatives differ only inside one \(\mathcal T_\Sigma\)-orbit, and the quotient-local
physical algebra \(\mathcal A_{\mathrm{phys}}(U)\) of Definition~\ref{def:physobs} is defined
precisely on that orbit space. Therefore the induced physical state and all expectations of
physical observables agree there as well.
\end{proof}


\section{Assumption-Dependent Distributed-Systems and QECC Extensions of the Consensus Formalism}\label{app:bft-qecc}
\paragraph{Support labels.}
\begin{description}
  \item[\textbf{[Established]}] Follows from cited prior work or a complete
        argument given here.
  \item[\textbf{[Assumption-dependent]}] True under additional assumptions not
        derived from \textrm{OPH} first principles.
  \item[\textbf{[Conjecture]}] A plausible open direction, not a settled result.
\end{description}

% ------------------------------------------------------------
\subsection*{B.1\quad Theorem 1: QBFT Safety Bound}
% ------------------------------------------------------------

\begin{definition}[QBFT-style protocol]
\label{def:qbft-style}
A consensus protocol is \emph{QBFT-style} in this analysis if it satisfies
the following three structural properties.
The safety proof of Theorem~\ref{thm:qbft-safety} uses all three;
the theorem does not hold for protocols lacking any of them without a
compensating change to the argument.
\begin{enumerate}[(P1)]
  \item \textbf{One-vote-per-view.}
        Each nonfaulty node casts at most one vote per view number.
        A node that has voted in view~$v$ ignores any later request
        to vote in view~$v$.
  \item \textbf{Certificate semantics.}
        A decision requires a valid quorum certificate: in the classical
        exact-size case used below, $q=2f+1$ distinct, unforgeable,
        authenticated votes for the same value in the same view.  For
        larger validator sets at fixed fault budget, the quorum threshold must
        be scaled so that the overlap condition in (A6) remains true.
  \item \textbf{DLS-style view-change.}
        If no certificate is produced within a timeout, every nonfaulty node
        increments the view number by one and a new leader is selected by a
        fixed deterministic rule.  At \textrm{GST}, timeouts fire correctly
        and the view-change terminates in bounded rounds.
\end{enumerate}
The Istanbul BFT / QBFT protocol family
\cite{Moniz2020,Saltini}
satisfies (P1)--(P3) and is the intended instance.
\end{definition}

\paragraph{Assumptions A1--A6.}
\begin{enumerate}[(A1)]
  \item \textbf{Partial synchrony (DLS).}
        Fixed but initially unknown bounds $\Delta$ (message delay) and
        $\Phi$ (processing rates).
        \emph{Safety} holds without extra timing assumptions;
        \emph{liveness} holds after the Global Stabilisation Time (\textrm{GST}).
  \item \textbf{Byzantine fault model.}
        At most $f$ observers behave arbitrarily; the remaining $n-f$ are nonfaulty.
  \item \textbf{Classical exact sizing for this fixed-quorum theorem.}
        $n = 3f+1$, so that $q=2f+1$ certificates have a nonfaulty overlap
        witness.  The usual resilience condition $n\ge 3f+1$ is necessary for
        the fault model, but when $n>3f+1$ a fixed $q=2f+1$ quorum is
        insufficient for the overlap used in the safety proof; one must either
        impose (A6) directly or choose a threshold $q$ with
        $2q\ge n+f+1$.
  \item \textbf{Strong quorum connectivity.}
        Every quorum $Q$ with $|Q|=q$ is strongly connected within~$G$:
        for any $u,v\in Q$ there is a directed path in~$G$ contained
        entirely in~$Q$.
        This is strictly stronger than requiring the overlap graph of
        quorums to be connected, and is needed to propagate signed votes
        within a quorum.
  \item \textbf{Message authentication.}
        All messages carry unforgeable digital signatures.
  \item \textbf{OPH quorum overlap.}
        Any two quorums $Q_a, Q_b$ of size $q$ satisfy
        $|Q_a\cap Q_b|\geq f+1$.  For $q=2f+1$ this is guaranteed by the
        exact sizing $n=3f+1$ in (A3); for general $n$ it is the separate
        threshold condition $2q\ge n+f+1$, not a consequence of
        $n\ge3f+1$ alone.
\end{enumerate}

D.\ Matscheko's Lean audit of this appendix checks the finite quorum-overlap
core and records the same boundary caveat: at fixed \(q=2f+1\), the overlap
step is exact-size \(n=3f+1\) logic unless (A6) is imposed separately.

\begin{theorem}[QBFT Safety Bound
  {\normalfont\small [Same-view case established under A1--A6;
  cross-view safety conditional on a (P4) locking rule]}]
\label{thm:qbft-safety}
Under assumptions (A1)--(A6), any consensus protocol satisfying
(P1)--(P3) of Definition~\ref{def:qbft-style} and run over the
\textrm{OPH} observer graph satisfies:
\begin{enumerate}[(i)]
  \item \textbf{Same-view safety.} No two nonfaulty observers finalise
        conflicting patch states in the same view.
  \item \textbf{Liveness.} After \textrm{GST}, every nonfaulty observer
        finalises within $O(f\cdot\Delta)$ wall-clock time.
  \item \textbf{Optimality.} The bound $f<n/3$ is tight.
\end{enumerate}
\end{theorem}

\paragraph{Cross-view boundary.}
Clause (i) is stated for same-view finalisation because the argument below
covers exactly that case. Cross-view safety in the IBFT/QBFT protocol family
uses an additional prepared-certificate locking property, an explicit
(P4) hypothesis under which a nonfaulty node votes in a later view only for
a value carried by the highest prepared certificate it has received.
(P1)--(P3) alone do not exclude finalisation of conflicting values across
views; a cross-view extension of this theorem requires (P4) and its
accompanying view-change justification argument.
On its stated assumptions, this finalisation theorem is a finite carrier of stack row C1 of
the program's consistency stack: record permanence, meaning no two nonfaulty observers finalise
conflicting patch states, enters as a consistency requirement on the observer net, not as an
added postulate.

\begin{proof}[Proof sketch]
\textit{Safety.}
Suppose $O_a$ and $O_b$ finalise $s_a\neq s_b$ in the same view.
By (P2), each required a certificate of $q$ votes: sets $Q_a,Q_b$.
By (A6), $|Q_a\cap Q_b|\geq f+1$.  In the classical exact-size case this is
the inclusion-exclusion calculation
$|Q_a\cap Q_b|\geq(2f+1)+(2f+1)-(3f+1)=f+1$.
By (A2), at most $f$ are Byzantine, so $Q_a\cap Q_b$ contains a
nonfaulty $O^*$.
By (A4), $O^*$'s signed vote is path-reachable within both quorums.
By (P1), $O^*$ voted for at most one value. Contradiction.

\textit{Liveness and Optimality}
follow from \cite{DLS88} (Thm.~4.4) and \cite{Lamport1982}, cited directly.

\textit{Note on FLP.}
Fischer, Lynch, Paterson~\cite{FLP85} is an impossibility result for
fully asynchronous systems; it does not bear on achievability under
partial synchrony (A1).
\end{proof}

% ------------------------------------------------------------
\subsection*{B.2\quad Theorem 2: Convergence of the OPH Repair Map}
% ------------------------------------------------------------

\begin{definition}[OPH Repair Map: Petz form]
\label{def:oph-repair-map}
Let $\sigma\in\mathcal{D}(\mathcal{H})$ be a full-rank reference state and
$\mathcal{N}:\mathcal{B}(\mathcal{H})\to\mathcal{B}(\mathcal{K})$ a quantum channel.
The \emph{OPH repair map} is
\[
  \mathcal{R}_{\sigma,\mathcal{N}}(\rho)
  := \sigma^{1/2}\,
     \mathcal{N}^\dagger\!\bigl(
       \mathcal{N}(\sigma)^{-1/2}\,\rho\,\mathcal{N}(\sigma)^{-1/2}
     \bigr)
     \,\sigma^{1/2},
\]
where $\mathcal{N}^\dagger$ is the adjoint channel and inverses are
taken on $\mathrm{supp}(\mathcal{N}(\sigma))$.
\end{definition}

\begin{remark}[Petz map vs.\ trace-distance projection]
The closest-point trace-distance projection
$\mathcal{P}_{\mathcal{S}}(\rho):=\arg\min_{\tau\in\mathcal{S}}\tfrac12\|\rho-\tau\|_1$
is a different object from the Petz map: it is defined by a variational
problem in trace-norm geometry and is not \textrm{CPTP} in general.
The two coincide only in very special cases not automatic in the
\textrm{OPH} setting.
All subsequent properties refer exclusively to
Definition~\ref{def:oph-repair-map}.
\end{remark}

\begin{proposition}[Petz map CPTP: domain-restricted statement
  {\normalfont\small [Established, subject to domain restriction]}]
\label{prop:petz-cptp}
Let $\sigma$ have full support on $\mathcal{H}$.
\begin{enumerate}[(a)]
  \item $\mathcal{R}_{\sigma,\mathcal{N}}$ is completely positive.
  \item $\mathcal{R}_{\sigma,\mathcal{N}}$ is trace-preserving on
        $\mathrm{supp}(\mathcal{N}(\sigma))$, i.e., on inputs $\rho$
        for which $\mathcal{N}(\sigma)^{-1/2}\rho\,\mathcal{N}(\sigma)^{-1/2}$
        is well-defined.
  \item If additionally $\mathcal{N}(\sigma)$ has full rank on $\mathcal{K}$,
        then $\mathcal{R}_{\sigma,\mathcal{N}}$ is \textrm{CPTP} on all
        of $\mathcal{B}(\mathcal{K})$.
\end{enumerate}
If $\mathcal{N}(\sigma)$ is \emph{not} full rank on $\mathcal{K}$,
then either (i)~the domain must be restricted to
$\mathrm{supp}(\mathcal{N}(\sigma))$, or
(ii)~pseudoinverses must replace the inverses (generalised Petz map;
cf.~\cite{Junge2018}), or
(iii)~a regularisation
$\mathcal{N}(\sigma)\mapsto\mathcal{N}(\sigma)+\varepsilon\mathbf{1}$
must be introduced.
Note that full-rank $\sigma$ does \emph{not} prevent $\mathcal{N}(\sigma)$
from being rank-deficient: the channel may map the support of $\sigma$
into a strict subspace of $\mathcal{K}$.
In the \textrm{OPH} setting, whether $\mathcal{N}(\sigma)$ is full rank
depends on the specific overlap channel and must be verified for the chosen analytic channel
model. The finite OPH repair theorem instead uses the declared Lyapunov descent law on a finite
patch net.
\end{proposition}

\begin{proof}
Complete positivity follows from composing three \textrm{CP} operations:
\begin{enumerate}[(i)]
  \item sandwiching by
        $\mathcal{N}(\sigma)^{-1/2}(\cdot)\mathcal{N}(\sigma)^{-1/2}$
        on $\mathrm{supp}(\mathcal{N}(\sigma))$;
  \item $\mathcal{N}^\dagger$;
  \item sandwiching by $\sigma^{1/2}(\cdot)\sigma^{1/2}$.
\end{enumerate}
Trace preservation in the full-rank case:
Petz~\cite{Petz1986}; Fagnola--Umanit\`{a}~\cite{FagnolaUmanita2010}.
\end{proof}

\begin{proposition}[Analytic contraction certificate {\normalfont\small [Assumption-dependent]}]
\label{prop:petz-contraction}
Suppose a declared quotient repair map \(T:Q\to Q\) on a metric physical quotient
\((Q,d_Q)\) is strictly contractive with coefficient \(\lambda\in(0,1)\):
\[
  d_Q(Tx,Ty)\le \lambda d_Q(x,y).
\]
Then \(T\) has at most one fixed point. If a fixed point \(x_\star\) exists, the ideal iterates
obey \(d_Q(T^t x,x_\star)\le \lambda^t d_Q(x,x_\star)\).
This is an analytic contraction condition. It is not required for the finite OPH normal-form
theorem, where termination follows from strict Lyapunov descent of accepted repairs.
\end{proposition}

\begin{theorem}[Noisy approximate repair stability {\normalfont\small [Contraction branch]}]
\label{thm:noisy-contraction}
Assume the contraction certificate of Proposition~\ref{prop:petz-contraction}, and let
\(\widetilde T:Q\to Q\) be an implemented noisy repair map satisfying
\[
d_Q(\widetilde T x,Tx)\le \varepsilon
\qquad\text{for all }x\in Q.
\]
If \(x_\star\) is the ideal fixed point of \(T\), then
\[
d_Q(\widetilde T^t x,x_\star)
\le
\lambda^t d_Q(x,x_\star)+\frac{\varepsilon}{1-\lambda}.
\]
Thus the noisy branch converges only to a controlled error ball, and only after the contraction
certificate and uniform implementation-error bound are supplied.
\end{theorem}

\begin{proof}
The recursion
\[
d_Q(\widetilde T x,x_\star)
\le d_Q(\widetilde T x,Tx)+d_Q(Tx,Tx_\star)
\le \varepsilon+\lambda d_Q(x,x_\star)
\]
iterates to the displayed geometric-series bound.
\end{proof}

\begin{proposition}[Spectral-gap criterion {\normalfont\small [Model-dependent]}]
\label{conj:spectral-gap}
Let \(\mathcal T\) be the Markov, channel, or transfer operator induced by iterated OPH repair on
a declared analytic realization. If \(\mathcal T\) has stationary projection \(\Pi_\star\) and
constants \(C<\infty\), \(\delta>0\) such that on the nonstationary subspace
\[
\|\mathcal T^t-\Pi_\star\|\le C e^{-\delta t},
\]
then the corresponding analytic channel model has exponential convergence. The finite OPH repair
package supplies termination by Lyapunov descent on its declared finite state space; a spectral
gap is a separate quantitative mixing condition for this BFT/QECC-style extension.
\end{proposition}

\begin{theorem}[Exponential Convergence
  {\normalfont\small [Under Proposition~\ref{conj:spectral-gap}]}]
\label{thm:exp-convergence}
Under Proposition~\ref{conj:spectral-gap}, for any initial analytic state \(\rho\),
\[
  \bigl\|\mathcal T^t\rho-\Pi_\star\rho\bigr\|
  \leq C\,e^{-\delta t}\bigl\|\rho-\Pi_\star\rho\bigr\|.
\]
This theorem belongs only to the spectral-gap branch. It is not a consequence of a bare finite
overlap graph or of finite Lyapunov descent alone.
\end{theorem}

% ------------------------------------------------------------
\subsection*{B.3\quad Theorem 3: QECC Correspondence}
% ------------------------------------------------------------

\paragraph{Notation.}
$N=\dim(\mathcal{H})=2^n$ for $n$ physical qubits.
Standard notation: $[[n,k,d]]$ stabilizer code; $K=2^k$; quantum
Singleton bound: $k\leq n-2(d-1)$.

\begin{theorem}[No free min-cut theorem for bare overlap graphs
  {\normalfont\small [Established]}]
\label{thm:no-free-mincut}
Let \(G=(V,E)\) be any connected graph with \(|V|\ge2\). The graph \(G\) alone does not determine
the Hamming distance of the consistency set \(C\) of a finite overlap net on \(G\). In particular,
the same graph can realize a binary constraint code of distance \(1\) or a binary repetition code
of distance \(|V|\).
\end{theorem}

\begin{proof}
Set \(S_i=\{0,1\}\) for every vertex. For the repetition realization, choose \(I_e=\{0,1\}\) and
let both endpoint readouts be the identity. Then every edge imposes \(x_i=x_j\), so the only
global codewords are \(00\cdots0\) and \(11\cdots1\), whose Hamming distance is \(|V|\).

For the trivial-overlap realization, keep the same vertex state spaces but let every endpoint
readout be the constant map to \(0\). Then every binary assignment is globally consistent, so the
minimum Hamming distance among distinct codewords is \(1\). The graph is unchanged. Therefore
distance is a property of the code realization--state spaces, readout maps, logical dictionary,
metric, and error model--not of the bare overlap graph.
\end{proof}

\begin{definition}[Topological-code realization certificate]\label{def:tcert}
An OPH overlap network may be treated as a QECC/topological code only after supplying a tuple
\[
\mathsf{TCert}=
(K,\mathcal H_{\mathrm{phys}},\mathcal H_{\mathrm{code}},
\partial_2,\partial_1,S_X,S_Z,\mathcal L_X,\mathcal L_Z,\mathcal E,\mathcal R),
\]
where \(C_2\xrightarrow{\partial_2}C_1\xrightarrow{\partial_1}C_0\) is a chain complex over
\(\mathbb F_2\), the physical carriers live in \(\mathcal H_{\mathrm{phys}}\), the protected
subspace is \(\mathcal H_{\mathrm{code}}\), \(S_X,S_Z\) are stabilizer or gauge checks,
\(\mathcal L_X,\mathcal L_Z\) are logical-operator classes, \(\mathcal E\) is a declared error
family, and \(\mathcal R\) is a recovery map or recovery family.
\end{definition}

\begin{theorem}[Certified topological-code distance and min-cut
  {\normalfont\small [Assumption-dependent]}]
\label{thm:tcert-distance}
Suppose a certificate \(\mathsf{TCert}\) of Definition~\ref{def:tcert} is supplied, with logical
classes identified as
\[
\mathcal L_X\simeq H_1(K;\mathbb F_2),
\qquad
\mathcal L_Z\simeq H^1(K;\mathbb F_2),
\]
and with boundary conditions excluding lower-weight trivial representatives. Then the certified
distance is
\[
d=
\min\left\{
\min_{\ell\in\mathcal L_X\setminus0}|\ell|,
\min_{\ell^\star\in\mathcal L_Z\setminus0}|\ell^\star|
\right\}.
\]
Only in geometries where this homological systole equals the relevant graph min-cut may one write
\(d=\mathrm{mincut}(G_{\mathrm{OPH}})\) \cite{Kitaev2003,Dennis2002}.
\end{theorem}

\begin{proof}
This is the standard stabilizer/topological-code distance statement once the chain complex,
checks, logical representatives, and boundary conditions are declared. The min-cut equality is an
additional geometric identification of that homological minimum with a graph cut. By
Theorem~\ref{thm:no-free-mincut}, it cannot be inferred from the bare graph.
\end{proof}

\begin{conjecture}[Communication complexity {\normalfont\small [Conjecture]}]
\label{conj:comm-complexity}
The \textrm{OPH} consensus-repair protocol, realised as a quantum
communication task, has per-round complexity $O(n\cdot\mathrm{poly}(d))$
for a chosen communication encoding (cf.~\cite{BCW98}). The fixed finite repair theorem gives
termination after a supplied descent law; it does not by itself fix a quantum communication
complexity class for every implementation.
\end{conjecture}

\begin{theorem}[QECC resilience under a supplied code certificate
  {\normalfont\small [Assumption-dependent]}]
\label{thm:qecc}
Assume a genuine code subspace \(\mathcal H_{\mathrm{code}}\subseteq\mathcal H_{\mathrm{phys}}\)
with projector \(\Pi\), and let \(\mathcal E_t\) be the declared set of errors supported on fewer
than \(t\) corrupted patches or physical carriers. If
\[
\Pi E_a^\dagger E_b \Pi=\alpha_{ab}\Pi
\qquad
\forall E_a,E_b\in\mathcal E_t,
\]
then there exists a recovery channel correcting all errors in \(\mathcal E_t\)
\cite{KL97}. If the supplied certificate also gives distance \(d\), then all errors of weight
\(t<d/2\) are correctable. No such resilience statement follows from the bare overlap graph.
\end{theorem}

\begin{corollary}[QECC extension inventory]\label{cor:qecc-inventory}
Under Definition~\ref{def:tcert}, Theorem~\ref{thm:tcert-distance}, and
Theorem~\ref{thm:qecc}, the OPH BFT/QECC extension carries the following claim split:
\begin{enumerate}[(i)]
  \item \textrm{[Assumption-dependent]} Code distance is the certified homological minimum; it
        equals \(\mathrm{mincut}(G_{\mathrm{OPH}})\) only under the additional systole/min-cut
        identification.
  \item \textrm{[Established]} The Knill--Laflamme \textrm{QECC}
        theorem supplies recovery once the projector and error family satisfy the displayed
        condition.
  \item \textrm{[Conjecture]} Per-round communication complexity is
        $O(n\cdot\mathrm{poly}(d))$.
\end{enumerate}
\end{corollary}

% ------------------------------------------------------------
\subsection*{B.4\quad Theorem 4: Asynchronous Convergence}
% ------------------------------------------------------------

\paragraph{Why fairness alone does not give a probability-1, spectral, or wall-clock statement.}
Standard strong fairness guarantees that every enabled action fires infinitely often along any fair
schedule; it does not impose a probability space on schedules, a transfer-operator spectral gap,
or a message-delay bound. A convergence statement of the form ``converges with probability~1''
requires a measure on schedules. Exponential convergence requires the spectral-gap certificate of
Theorem~\ref{thm:exp-convergence}. Bounded wall-clock liveness requires partial synchrony and
quorum assumptions. The \textrm{FLP} impossibility result~\cite{FLP85} confirms that fairness is
insufficient for bounded-time consensus in a fully asynchronous system.

\paragraph{Additional assumptions for a quantitative bound.}
\begin{enumerate}[(B1)]
  \item Finite known bound $\Delta$ on message delay after \textrm{GST}.
  \item Finite bound $\Phi$ on processing rates.
  \item $f < n/3$.
\end{enumerate}

\begin{theorem}[Eventual finite repair termination
  {\normalfont\small [Finite-descent branch]}]
\label{thm:eventual-convergence}
For a finite OPH patch net with a strict Lyapunov-decreasing accepted repair relation, every
maximal repair run terminates after finitely many accepted repairs. The step count is bounded by
the finite value-set bound of Proposition~\ref{prop:finite-step-bound}. If the local-diamond and
repair-completeness clauses also hold, Newman's lemma upgrades this termination statement to the
unique schedule-independent quotient normal form of Theorem~\ref{thm:confluence}.

This theorem is finite descent convergence in repair steps. It is not trace-norm convergence of a
Petz channel, not probability-one convergence over random schedules, not exponential convergence,
and not a wall-clock liveness theorem.
\end{theorem}

\begin{theorem}[Quantitative Convergence
  {\normalfont\small [Assuming (B1)--(B3)]}]
\label{thm:quantitative-convergence}
In a partially synchronous \textrm{OPH} observer network satisfying
(B1)--(B3), after \textrm{GST} every nonfaulty observer reaches consensus
within $T=O(f\cdot\Delta)$ wall-clock time
(by applying the \textrm{DLS} framework~\cite{DLS88}, Thm.~4.4,
to the \textrm{OPH} repair protocol;
requires (B1) and (B2) explicitly and does not follow from fairness alone).
\end{theorem}

% ------------------------------------------------------------
\subsection*{B.5\quad Extension Boundaries}
% ------------------------------------------------------------

\begin{itemize}
  \item A bare OPH overlap graph is a finite constraint-code presentation only. Its graph does
        not determine code distance, correctable error weight, or min-cut resilience.
  \item Analytic spectral-gap and full-rank estimates for a chosen stochastic or channel-level
        BFT/QECC realization are model-specific refinements. They are separate from the finite
        OPH repair theorem, where the accepted repair law supplies Lyapunov descent directly.
  \item Long-run noisy approximate consensus is available only on the fair-block contraction
        branch of Theorem~\ref{thm:noisy-fair-block-consensus}: the chosen implementation must
        certify fair blocks, expected contraction toward the exact quotient normal-form set, and
        controlled within-block excursions. Fairness alone does not supply that certificate.
  \item Topological-code distance equals graph min-cut only after a concrete chain complex,
        boundary condition, logical-operator dictionary, error family, and systole/min-cut
        identification are supplied for the chosen code realization.
  \item Communication-complexity bounds require a concrete quantum communication encoding and
        implementation cost model. They are not consequences of finite normal-form termination
        alone.
  \item The core OPH consensus paper supplies observable-level confluence and refinement-limit
        normal-form/holonomy classes on their declared theorem surfaces; the BFT/QECC statements
        above are separate protocol-style extensions.
\end{itemize}

% ------------------------------------------------------------
%  Bibliography entries for this appendix
%  (merge into main .bib file or \bibliography)
% ------------------------------------------------------------
%
% @article{FLP85,
%   author  = {Fischer, M.J. and Lynch, N.A. and Paterson, M.S.},
%   title   = {Impossibility of Distributed Consensus with One Faulty Process},
%   journal = {Journal of the ACM},
%   volume  = {32}, number = {2}, pages = {374--382}, year = {1985}}
%
% @article{Lamport1982,
%   author  = {Lamport, L. and Shostak, R. and Pease, M.},
%   title   = {The {Byzantine} Generals Problem},
%   journal = {ACM Transactions on Programming Languages and Systems},
%   volume  = {4}, number = {3}, pages = {382--401}, year = {1982}}
%
% @article{DLS88,
%   author  = {Dwork, C. and Lynch, N.A. and Stockmeyer, L.},
%   title   = {Consensus in the Presence of Partial Synchrony},
%   journal = {Journal of the ACM},
%   volume  = {35}, number = {2}, pages = {288--323}, year = {1988}}
%
% @inproceedings{Castro1999,
%   author    = {Castro, M. and Liskov, B.},
%   title     = {Practical {Byzantine} Fault Tolerance},
%   booktitle = {OSDI 1999}, pages = {173--186}, year = {1999}}
%
% @misc{Moniz2020,
%   author = {Moniz, H.},
%   title  = {The {Istanbul BFT} Consensus Algorithm},
%   eprint = {2002.03613}, year = {2020}}
%
% @misc{Saltini,
%   author = {Saltini, R. and others},
%   title  = {{QBFT} Formal Specification and Verification},
%   howpublished = {\url{https://github.com/Consensys/qbft-formal-spec-and-verification}}}
%
% @article{Petz1986,
%   author  = {Petz, D.},
%   title   = {Sufficient subalgebras and the relative entropy of states
%              of a von Neumann algebra},
%   journal = {Communications in Mathematical Physics},
%   volume  = {105}, number = {1}, pages = {123--131}, year = {1986}}
%
% @article{FagnolaUmanita2010,
%   author  = {Fagnola, F. and Umanit\`{a}, V.},
%   title   = {Generators of detailed balance quantum {Markov} semigroups},
%   journal = {Infinite Dimensional Analysis, Quantum Probability},
%   volume  = {13}, number = {3}, pages = {459--486}, year = {2010}}
%
% @article{Junge2018,
%   author  = {Junge, M. and others},
%   title   = {Universal recovery maps and approximate sufficiency of
%              quantum relative entropies},
%   journal = {Annales Henri Poincar\'e},
%   volume  = {19}, number = {8}, pages = {2505--2555}, year = {2018}}
%
% @article{KL97,
%   author  = {Knill, E. and Laflamme, R.},
%   title   = {Theory of quantum error-correcting codes},
%   journal = {Physical Review A},
%   volume  = {55}, number = {2}, pages = {900--911}, year = {1997}}
%
% @phdthesis{Gottesman1997,
%   author = {Gottesman, D.},
%   title  = {Stabilizer Codes and Quantum Error Correction},
%   school = {California Institute of Technology}, year = {1997}}
%
% @article{Kitaev2003,
%   author  = {Kitaev, A.},
%   title   = {Fault-tolerant quantum computation by anyons},
%   journal = {Annals of Physics},
%   volume  = {303}, number = {1}, pages = {2--30}, year = {2003}}
%
% @article{Dennis2002,
%   author  = {Dennis, E. and Kitaev, A. and Landahl, A. and Preskill, J.},
%   title   = {Topological quantum memory},
%   journal = {Journal of Mathematical Physics},
%   volume  = {43}, number = {9}, pages = {4452--4505}, year = {2002}}
%
% @inproceedings{BCW98,
%   author    = {Buhrman, H. and Cleve, R. and Wigderson, A.},
%   title     = {Quantum vs.\ Classical Communication and Computation},
%   booktitle = {STOC 1998}, pages = {63--68}, year = {1998}}
%
% @book{Attiya2004,
%   author    = {Attiya, H. and Welch, J.},
%   title     = {Distributed Computing: Fundamentals, Simulations, and
%                Advanced Topics},
%   publisher = {Wiley-Interscience}, year = {2004}}


\begin{thebibliography}{99}

\bibitem{OPH}
B.~M\"uller, A.~Osika, M.~Poneder, K.~Xue, B.~Cassie, P.~Nguyen, M.~A.~Visser, and K.~A.~Anirudha,
\emph{Observers Are All You Need},
2026.\\
Available at \url{https://oph-book.floatingpragma.io/papers/observers_are_all_you_need.pdf}.

\bibitem{newman42}
M.~H.~A.~Newman,
``On theories with a combinatorial definition of `equivalence',''
\emph{Ann.\ of Math.} \textbf{43} (1942), no.~2, 223--243.

\bibitem{FLP85}
M.~J.~Fischer, N.~A.~Lynch, and M.~S.~Paterson,
``Impossibility of distributed consensus with one faulty process,''
\emph{J.\ ACM} \textbf{32} (1985), no.~2, 374--382.

\bibitem{Lamport1982}
L.~Lamport, R.~Shostak, and M.~Pease,
``The Byzantine generals problem,''
\emph{ACM Trans.\ Program.\ Lang.\ Syst.} \textbf{4} (1982), no.~3, 382--401.

\bibitem{DLS88}
C.~Dwork, N.~A.~Lynch, and L.~Stockmeyer,
``Consensus in the presence of partial synchrony,''
\emph{J.\ ACM} \textbf{35} (1988), no.~2, 288--323.

\bibitem{Moniz2020}
H.~Moniz,
\emph{The Istanbul BFT Consensus Algorithm},
arXiv:2002.03613, 2020.

\bibitem{Saltini}
R.~Saltini et al.,
\emph{QBFT Formal Specification and Verification}.
Available at \url{https://github.com/Consensys/qbft-formal-spec-and-verification}.

\bibitem{Petz1986}
D.~Petz,
``Sufficient subalgebras and the relative entropy of states of a von Neumann algebra,''
\emph{Commun.\ Math.\ Phys.} \textbf{105} (1986), no.~1, 123--131.

\bibitem{FagnolaUmanita2010}
F.~Fagnola and V.~Umanit\`a,
``Generators of detailed balance quantum Markov semigroups,''
\emph{Infinite Dimensional Analysis, Quantum Probability and Related Topics}
\textbf{13} (2010), no.~3, 459--486.

\bibitem{Junge2018}
M.~Junge et al.,
``Universal recovery maps and approximate sufficiency of quantum relative entropies,''
\emph{Ann.\ Henri Poincar\'e} \textbf{19} (2018), no.~8, 2505--2555.

\bibitem{KL97}
E.~Knill and R.~Laflamme,
``Theory of quantum error-correcting codes,''
\emph{Phys.\ Rev.\ A} \textbf{55} (1997), no.~2, 900--911.

\bibitem{Kitaev2003}
A.~Kitaev,
``Fault-tolerant quantum computation by anyons,''
\emph{Ann.\ Phys.} \textbf{303} (2003), no.~1, 2--30.

\bibitem{Dennis2002}
E.~Dennis, A.~Kitaev, A.~Landahl, and J.~Preskill,
``Topological quantum memory,''
\emph{J.\ Math.\ Phys.} \textbf{43} (2002), no.~9, 4452--4505.

\bibitem{BCW98}
H.~Buhrman, R.~Cleve, and A.~Wigderson,
``Quantum vs.\ classical communication and computation,''
in \emph{Proceedings of STOC 1998}, pp.~63--68.

\bibitem{fawziRenner15}
O.~Fawzi and R.~Renner,
``Quantum conditional mutual information and approximate Markov chains,''
\emph{Commun.\ Math.\ Phys.} \textbf{340} (2015), 575--611,
arXiv:1410.0664.

\end{thebibliography}

\end{document}
